\section{Integral and $\lambda$-ring controls on the bridges}\label{sec:root-findings}
The results in this section are deductions checked in this run. They are not claims of historical priority. We distinguish an identity of characters, a decomposition in a Grothendieck group, an actual direct-sum decomposition, and a dimension identity.

\begin{proposition}[G-01: The three meanings of the second operation]\label{G-01}
Let $E$ be a finite $G$-set. Write $b_2(E)$ for the $G$-set of two-element subsets and $s_2(E)$ for the $G$-set of two-element multisets. In the Burnside ring one has
\[
 [s_2(E)]=[b_2(E)]+[E],\qquad [E]^2=[E]+[\operatorname{Inj}(\{1,2\},E)].
\]
The identity $[E]^2=2[b_2(E)]+[E]$ does not hold in general. If $V$ is a representation over a field of characteristic zero, the correct special $\lambda$-ring identities are
\[
 [\operatorname{Sym}^2V]=[V]^2-[\wedge^2V],\quad
 [\operatorname{Sym}^2V]-[\wedge^2V]=\psi^2[V],\quad
 [V]^2=2[\wedge^2V]+\psi^2[V].
\]
Consequently the expression $2\lambda^2(E)+E=E^2$ in cote 133, scan 019, line 53 is valid after taking cardinalities, but is not a universal equivariant or special-$\lambda$-ring identity. The supplied study's assertion that the two equivariant descriptions agree needs this correction.
\end{proposition}
\begin{proof}
Partition a multiset of cardinality two into the diagonal multisets $\{e,e\}$ and the subsets $\{e,e'\}$ with $e\ne e'$. This partition is $G$-equivariant. Similarly partition ordered pairs into diagonal and off-diagonal pairs. For a counterexample take $G=C_2$ and its regular two-element set $E$. Then $E\times E$ is a disjoint union of two free $C_2$-orbits, whereas $b_2(E)$ is one fixed point. Thus the asserted equality would read $2[E]=[E]+2[\mathrm{pt}]$. Applying the fixed-point homomorphism for the whole group gives $0=2$, a contradiction.
For representations, first prove the universal identity on $\mathrm{GL}_r$ and then pull it back along the representation of $G$. Write the character of its defining representation on the diagonal torus as $x_1+\cdots+x_r$. The symmetric and exterior squares have characters $\sum_{i\leq j}x_ix_j$ and $\sum_{i<j}x_ix_j$. Their sum is $(\sum_i x_i)^2$, and their difference is $\sum_i x_i^2=\psi^2[V]$. Character injectivity in characteristic zero proves the formulas. In the same $C_2$ example the permutation representation is $1\oplus\epsilon$, and $\wedge^2V=\epsilon$, whereas linearization of $b_2(E)$ is $1$. This explicitly shows why the subset operation cannot be silently identified with exterior power. Taking dimensions sends $\psi^2[V]$ to $\dim V$ and recovers the numerical identity. None of this determines the ink on a scan which has not been supplied.
\end{proof}

\begin{proposition}[G-02: The terminal exterior-power variable is necessary]\label{G-02}
For $i\geq2$, a universal formula for $\lambda^i(xy)$ in special $\lambda$-rings cannot in general omit $\lambda^i(x)$ or $\lambda^i(y)$. Thus the terminal index in cote 133, scan 006, lines 25--26 must be $i$ if those lines state the universal product formula.
\end{proposition}
\begin{proof}
Set $y=1$, for which $\lambda^1(1)=1$ and $\lambda^j(1)=0$ for $j>1$. A formula using only $x,\lambda^2(x),\ldots,\lambda^{i-1}(x)$ would express $\lambda^i(x)$ in the subring generated by these earlier operations. Evaluate at the defining representation of $\mathrm{GL}_i$ and restrict to its diagonal torus. This would express the elementary symmetric polynomial $e_i$ in $\mathbb Z[e_1,\ldots,e_{i-1}]$, contrary to the algebraic independence of $e_1,\ldots,e_i$ (equivalently the fundamental theorem of symmetric polynomials). Interchange $x$ and $y$ for the other endpoint. The argument establishes the required mathematical reading, not a new visual transcription.
\end{proof}

\begin{proposition}[G-03: Finite $\gamma$-data do not determine multiplicities]\label{G-03}
For every integer $D\geq0$ there are two honest representations $A_D,B_D$ of the split group scheme $\mathbb G_{m,\mathbb Z}$, each of rank $2^D$, which have the same class modulo $F^{D+1}_\gamma$ but different weight multiplicities. Over $\mathbb Q$ their Chern characters agree in degrees $0,\ldots,D$.
\end{proposition}
\begin{proof}
Let $L$ denote the standard character and put $r=D+1$. Define
\[
 A_D=\bigoplus_{\substack{0\leq j\leq r\\j\text{ even}}}L^{\otimes j\,\oplus\binom rj},\qquad
 B_D=\bigoplus_{\substack{0\leq j\leq r\\j\text{ odd}}}L^{\otimes j\,\oplus\binom rj}.
\]
The two binomial parity sums are $2^{r-1}$. In the representation ring $\mathbb Z[L,L^{-1}]$ the difference of the classes is $(1-L)^r$, a product of $r$ rank-zero elements and hence an element of $F^r_\gamma$. The constant weight has multiplicity one in $A_D$ and zero in $B_D$, so the characters differ. With $c=c_1(L)$, the difference of the Chern characters is $(1-e^c)^r$, whose lowest possible degree is $r$. This proves the assertion. For a line bundle on a scheme, applying the associated-bundle map gives the identical statement in its $K_0$. This does not assert that the two resulting bundles are nonisomorphic on every scheme.
\end{proof}

\begin{proposition}[G-04: An integral coefficient model and its exact specialization gate]\label{G-04}
Let $W_{\mathbb Z}$ be a finite free module, let $\rho:\mathrm{GL}_{N,\mathbb Z}\to\mathrm{GL}(W_{\mathbb Z})$ be a homogeneous polynomial representation of positive degree $e$, and let $f\in W_{\mathbb Z}$. Define $X_{\mathbb Z}$ as the affine schematic closure of its orbit. Its coordinate algebra is torsion-free, hence flat over $\mathbb Z$. In degree $d$ its orbit-coefficient map is an explicit finite integer matrix
\[
 A_d:\operatorname{Sym}^d(W_{\mathbb Z}^*)\longrightarrow
 \mathbb Z[g_{11},\ldots,g_{NN}]_{\leq ed},\qquad P\longmapsto P(\rho(g)f).
\]
Its rank over $\mathbb Q$ is the degree-$d$ coordinate-ring dimension. If $\delta_d$ is any nonzero maximal-rank minor of $A_d$, then this dimension and the degree-$d$ kernel commute with base change over $\mathbb Z[1/\delta_d]$. At a prime dividing $\delta_d$, rank must be tested anew; the fiber of the integral closure need not be the schematic orbit closure of the reduced point.
\end{proposition}
\begin{proof}
The coordinate homomorphism of the orbit map has target $\mathbb Z[\mathrm{GL}_N]=\mathbb Z[g_{ij},\det(g)^{-1}]$, a torsion-free $\mathbb Z$-module. Its kernel $I$ defines the schematic closure and its image $\mathbb Z[W]/I$ is torsion-free. A module over the principal ideal domain $\mathbb Z$ is flat if it is torsion-free, proving the first assertion. Since the representation is polynomial, each degree-$d$ homogeneous coordinate polynomial maps to a polynomial of degree at most $ed$ in the entries of $g$. Equality on invertible matrices is equality in the polynomial ring because localization by the nonzero polynomial $\det(g)$ is injective. This proves the finite coefficient-matrix description.
Let $r=\operatorname{rank}_{\mathbb Q}A_d$. All $(r+1)$-minors vanish as integers, while an $r$-minor $\delta_d$ is nonzero. Over $\mathbb Z[1/\delta_d]$, use that invertible block for row and column elimination. The vanishing of larger minors makes the remaining Schur complement zero, so the matrix is equivalent to $\operatorname{diag}(I_r,0)$. Its image and kernel are direct summands there and commute with every base change. Without inversion, an inclusion of lattices need not remain injective modulo a prime: already multiplication by $p$ on $\mathbb Z$ has rank one over $\mathbb Q$ and rank zero modulo $p$. Thus no unsupported all-primes orbit-fiber conclusion follows. Homogeneity sends degree $d$ to degree $ed$, so the kernel is homogeneous and the stated grading is well defined. If $r=0$, take $\delta_d=1$.
\end{proof}

\begin{proposition}[G-05: Ternary saturation and the three-row exclusion]\label{G-05}
Work over $\mathbb C$, and let $n\geq2$, $N=n^2$, and $\mathrm{Det}_n\subseteq\operatorname{Sym}^n(\mathbb C^N)^*$ be the determinant orbit closure. Then $\mathrm{Det}_n$ contains every degree-$n$ form supported on a fixed three-dimensional subspace. Consequently, for every partition $\lambda\vdash nd$ with $\ell(\lambda)\leq3$,
\[
 \operatorname{mult}_{S_\lambda\mathbb C^N}\mathbb C[\mathrm{Det}_n]_d
 =\operatorname{mult}_{S_\lambda\mathbb C^N}\operatorname{Sym}^d(\operatorname{Sym}^n\mathbb C^N).
\]
No multiplicity obstruction with at most three rows can separate any other orbit closure of a degree-$n$ form in the same ambient space, including a padded permanent, from $\mathrm{Det}_n$.
\end{proposition}
\begin{proof}
We use the established linear determinantal representation theorem for smooth plane curves: a smooth plane degree-$n$ curve together with a line bundle $L$ of degree $g-1$ with $H^0(L)=0$ admits an $n\times n$ matrix of linear forms whose determinant is its defining form up to a nonzero scalar (Beauville, 2000, Proposition 3.1(a)). Such an $L$ exists over $\mathbb C$: for $g\geq1$, the effective locus in $\operatorname{Pic}^{g-1}$ is the image of the symmetric power $C^{(g-1)}$, of dimension at most $g-1$, whereas $\operatorname{Pic}^{g-1}$ has dimension $g$; for $g=0$, take the degree-$-1$ line bundle. Scaling one row absorbs the nonzero scalar. Hence every smooth ternary degree-$n$ form is obtained by a linear substitution in $\det_n$. Any endomorphism substitution is a limit of invertible substitutions because $\mathrm{GL}_N$ is Zariski dense in $\operatorname{End}(\mathbb C^N)$; continuity of the polynomial action therefore puts these forms in $\mathrm{Det}_n$. Smooth ternary forms form a dense open set, so its closedness puts every ternary form there.
Let $W=\operatorname{Sym}^d(\operatorname{Sym}^n\mathbb C^N)$, and choose the Borel convention whose highest-weight vector has nonnegative weight $\lambda=(\lambda_1,\lambda_2,\lambda_3,0,\ldots,0)$. In a monomial basis of $W$, zero weight in positions $4,\ldots,N$ forces every coefficient variable to involve only the first three basis vectors. Thus the entire highest-weight space $W_\lambda^{\mathrm{hw}}$ is contained in $\operatorname{Sym}^d(\operatorname{Sym}^n\mathbb C^3)$. A nonzero polynomial in these coefficient variables cannot vanish on all ternary forms. By the first paragraph, the degree-$d$ determinant ideal therefore meets $W_\lambda^{\mathrm{hw}}$ trivially. Exactness and complete reducibility in characteristic zero give the asserted multiplicity equality. Every other closed subvariety has a degree-$d$ coordinate ring which is a quotient of $W$, so its multiplicity cannot exceed this common ambient value. This proves the exclusion. The assertion is about complex characteristic-zero multiplicities, not an unqualified statement in positive characteristic.
\end{proof}

\begin{proposition}[G-06: The four-variable gate is nonvacuous]\label{G-06}
For $n\geq4$, the closure $\mathcal D_{n,4}$ of degree-$n$ quaternary forms which are determinants of $n\times n$ matrices of linear forms has dimension at most $2n^2+2$. Since
\[
 \binom{n+3}{3}-(2n^2+2)=\binom{n-1}{3}>0,
\]
this is a proper subvariety of the space of quaternary forms. The determinant ideal in $\operatorname{Sym}(\operatorname{Sym}^n\mathbb C^{n^2})$ therefore has a nonzero isotypic component with a partition of exactly four rows in some finite degree.
\end{proposition}
\begin{proof}
Consider the polynomial map
\[
 (A_1,A_2,A_3,A_4)\longmapsto\det(x_1A_1+x_2A_2+x_3A_3+x_4A_4)
\]
from an affine space of dimension $4n^2$. The group
$H=\{(P,Q)\in\mathrm{GL}_n^2:\det(P)=\det(Q)\}$ acts by $A_i\mapsto PA_iQ^{-1}$, preserving this map. It has dimension $2n^2-1$. Its generic stabilizer has dimension one. Indeed, for a tuple with $A_1=I$, a diagonal $A_2$ with pairwise distinct eigenvalues, and an $A_3$ with every off-diagonal entry nonzero, a stabilizer has $Q=P$; commuting with $A_2$ makes $P$ diagonal, and commuting with $A_3$ makes it scalar. The scalar pairs $(cI,cI)$ always stabilize every tuple. Upper semicontinuity of stabilizer dimension, or the rank calculation for the commutation equations on this exhibited open set, gives generic dimension one. Thus generic fibers contain orbits of dimension $2n^2-2$, and the dimension-of-fibers theorem yields the stated image bound. The displayed identity follows by expanding the binomial coefficient.
Let $\pi$ set all but the first four variables to zero. For every form in $\mathrm{Det}_n$, its image under $\pi$ lies in $\mathcal D_{n,4}$: approximate the form by invertible substitutions of $\det_n$, apply $\pi$, and use the closedness of $\mathcal D_{n,4}$. Choose a nonzero homogeneous equation $P$ of this proper cone. Then $P\circ\pi$ is a nonzero homogeneous equation of $\mathrm{Det}_n$, using only first-four coefficient variables. The ideal of $\mathcal D_{n,4}$ is $\mathrm{GL}_4$-stable; choose a highest-weight vector in a nonzero irreducible submodule of its degree-$d$ part. Embedded in the full coordinate ring, it is also a $\mathrm{GL}_{n^2}$ highest-weight vector: raising operators involving a variable beyond the fourth annihilate it. Its highest weight has at most four rows, and G-05 rules out at most three. Hence its length is exactly four. Existence of such an equation does not establish that it is nonzero on the padded permanent.
\end{proof}

\begin{proposition}[G-07: An exact four-row determinant-closure formula]\label{G-07}
Set $D=320112$. For every $d\geq0$ and partition $\lambda\vdash4d$ with $\ell(\lambda)\leq4$, padded by zeros to length four, the degree-$d$ determinant-closure multiplicity satisfies
\[
 m_{\mathrm{Det}_4}(\lambda,d)
 =a_\lambda(d[4])-a_{\lambda-(D,D,D,D)}((d-D)[4]).
\]
The second term is defined to be zero if $d<D$ or $\lambda_4<D$. In particular the four-row sector is ambient for $d<D$, and at $d=D$ the multiplicity of $\lambda=(D^4)$ is exactly one less than its ambient plethysm multiplicity.
\end{proposition}
\begin{proof}
Leal, Lozano Huerta, and Vite (2024), Theorem 2 and Table 2, identify the closure of linear determinantal quartic surfaces as an irreducible reduced hypersurface of degree $D$ in $\mathbb P(\operatorname{Sym}^4\mathbb C^4)^*$. Table 2 identifies this component $F_1$ by $a=(5,5,5,5)$, $b=(6,6,6,6)$, so every matrix entry has degree one. Its affine cone has a prime principal ideal $(F)$ in the polynomial ring $A=\operatorname{Sym}(\operatorname{Sym}^4\mathbb C^4)$, with $F$ homogeneous of degree $D$.
The line $\mathbb CF$ is stable under $\mathrm{GL}_4$. Every algebraic one-dimensional character of $\mathrm{GL}_4$ is a power of the determinant. Since scalar $zI$ acts on $A_D$ by $z^{4D}$, this character is $\det^D$. Thus multiplication by $F$ gives the exact sequence of $\mathrm{GL}_4$-modules
\[
 0\longrightarrow \det^D\otimes A_{d-D}\xrightarrow{\ \cdot F\ }A_d
 \longrightarrow(A/(F))_d\longrightarrow0.
\]
For $d<D$ its first term is zero.
Let $\pi$ restrict forms in sixteen variables to their first four variables. The image $\pi(\mathrm{Det}_4)$ is precisely the affine cone just described. One inclusion follows by continuity from the orbit of $\det_4$. For the converse, every four-variable linear determinant is an endomorphism substitution in $\det_4$, hence belongs to its orbit closure; the same is true of every limit of such forms, and such a form is fixed by $\pi$. It follows that the intersection of the full determinant ideal with the first-four coefficient polynomial ring is exactly $(F)$. Highest-weight spaces for partitions of at most four rows use only these coefficient variables, by the weight argument of G-05. Applying complete reducibility to the exact sequence and shifting a highest weight by the determinant character proves the formula. At $d=D$ the source is $\det^D\otimes\mathbb C$, so the deficit is one.
\end{proof}
\paragraph{Source correction for G-07.}
The attached Landsberg source, section \texttt{zhanglee}, theorem \texttt{zlthm}, records $640224$. The later result used here is the degree of the reduced prime divisor. Leal--Lozano Huerta--Vite, section 4.1 immediately before Theorem 2, explicitly gives $P_{20,2}=\tfrac12D_{23,14}$. This is the precise cycle-multiplicity distinction; using the unreduced Noether--Lefschetz coefficient as the polynomial degree doubles the answer. This correction concerns the cited literature, not a purported visual correction of cote 133.

\begin{proposition}[G-08: A uniform low-variable determinantal construction]\label{G-08}
Let $m\geq1$ and $k\geq2$, put $a=\lfloor k/2\rfloor$, $b=\lceil k/2\rceil$, and
\[
 B(m,k)=\binom{m+a}{a}+\binom{m+b}{b}-1.
\]
Over any commutative ring, every degree-$m$ homogeneous polynomial $q$ in $k$ variables admits a matrix of linear forms of size $B(m,k)$ with determinant
\[
 \ell^{B(m,k)-m}q,
\]
where $\ell$ may first be an independent variable and then be specialized to any linear form. The construction is polynomial in the coefficients of $q$ and defined over $\mathbb Z$.
Consequently, over $\mathbb C$, whenever $n\geq B(m,k)$, no multiplicity obstruction of length at most $k$ can separate the orbit closure of an arbitrarily padded degree-$m$ form $\ell^{n-m}f$ from $\mathrm{Det}_n$.
\end{proposition}
\begin{proof}
Split the variables into disjoint lists $x=(x_1,\ldots,x_a)$ and $y=(y_1,\ldots,y_b)$, and write
\[
 q=\sum_{|\alpha|+|\beta|=m}c_{\alpha\beta}x^\alpha y^\beta.
\]
Make a rooted directed tree with one vertex for every $x$-monomial of degree at most $m$. The parent of a nonconstant monomial is obtained by deleting one occurrence of its largest-index variable; orient the edge away from the constant root and label it by that variable. There is exactly one root-to-vertex path, with label product $x^\alpha$. Make the corresponding $y$-tree with every edge reversed, so the unique path from vertex $\beta$ to its constant root has product $y^\beta$. Connect vertex $\alpha$ in the first tree to vertex $\beta$ in the second by an edge of constant label $c_{\alpha\beta}$ exactly when $|\alpha|+|\beta|=m$. Every source-to-sink path uses exactly one connecting edge; its label product is $c_{\alpha\beta}x^\alpha y^\beta$. Therefore its path polynomial is exactly $q$. The graph is acyclic and has
\[
 s=\binom{m+a}{a}+\binom{m+b}{b}=B(m,k)+1
\]
vertices.
Choose a topological ordering with source first and sink last, and let $A$ be its strictly upper triangular weighted adjacency matrix. Then $\det(I-A)=1$ and $(I-A)^{-1}=I+A+\cdots+A^{s-1}$. The $(1,s)$ entry is the source-to-sink path polynomial $q$. The adjugate identity therefore expresses $q$, up to the fixed sign $(-1)^{s+1}$, as the determinant of the $(s-1)\times(s-1)$ minor of $I-A$ obtained by deleting row $s$ and column $1$. Change the sign of one row if needed. Its entries are affine linear in $x,y$.
Replace each constant entry $c$ of this minor by $c\ell$. Its determinant is the degree-$B(m,k)$ homogenization of $q$, hence is $\ell^{B(m,k)-m}q$. This identity is universal over $\mathbb Z[c_{\alpha\beta},x,y,\ell]$. Padding the matrix with $n-B(m,k)$ diagonal entries $\ell$ proves the corresponding size-$n$ statement.
For the multiplicity assertion, restrict any invertible substitution of $\ell^{n-m}f$ to $k$ coordinate variables. The result has the form $L^{n-m}q$, with $q$ a homogeneous degree-$m$ form in these variables. The construction puts this result in $\mathrm{Det}_n$. Continuity gives the same inclusion for the entire padded orbit closure. If $P$ is any highest-weight polynomial of length at most $k$ vanishing on $\mathrm{Det}_n$, it uses only coefficient variables supported on the first $k$ underlying variables (in the sense made precise in G-05), and therefore vanishes on the other orbit closure as well. Thus the highest-weight space of the determinant ideal embeds in that of the padded-form ideal. Subtracting their dimensions from the common ambient highest-weight multiplicity gives $m_{\mathrm{padded}}(\lambda,d)\leq m_{\mathrm{Det}_n}(\lambda,d)$ in every degree.
\end{proof}
\paragraph{Quantified implication.}
For $k=4$, the construction gives $B(m,4)=(m+1)(m+2)-1$. More generally $B(m,k)=O_k(m^{\lceil k/2\rceil})$. Thus for a fixed number of rows this sector is excluded once the padding exceeds the displayed bound. This is an explicit restricted lower-bound obstruction theorem, not a general separation result and not an efficient solution of arbitrary plethysm.

\begin{proposition}[G-09: A sharper low-row exclusion by shared exponent chains]\label{G-09}
Let $m\geq1$ and $k\geq3$. Define the following integer bounds. If $k=2r+1$ is odd, set
\[
 C(m,k)=2\binom{m+r}{r}-1.
\]
If $k=2r+2\geq4$ is even, for $1\leq h\leq m+1$ put $Q=\lfloor m/h\rfloor$, $s_0=\min(h-1,m)$ and
\[
\begin{aligned}
 C_h(m,k)&=\sum_{q=0}^{Q}\binom{m-qh+r}{r}+Q(h-1)
          +\sum_{s=0}^{s_0}\binom{m-s+r}{r}-1,\\
 C(m,k)&=\min_{1\leq h\leq m+1}C_h(m,k).
\end{aligned}
\]
Over a field, for every nonzero linear form $L$ and homogeneous degree-$m$ form $q$ in $k$ variables, $L^{C(m,k)-m}q$ has a linear determinantal representation of size $C(m,k)$. In coordinates with $L=x_0$ the construction is defined over $\mathbb Z$ and polynomial in the coefficients of $q$. For fixed $k\geq3$,
\[
 C(m,k)=O_k\bigl(m^{(k-1)/2}\bigr).
\]
Consequently, over $\mathbb C$, if $n\geq C(m,k)$, there is no multiplicity obstruction with at most $k$ rows separating any padded degree-$m$ form from $\mathrm{Det}_n$. This strengthens G-08 in the specified range.
\end{proposition}
\paragraph{Prior-result boundary.}
The asymptotic determinantal-representation order used here is established: Boralevi, van Doornmalen, Draisma, Hochstenbach, and Plestenjak (2017), Theorem 1.3, prove order $m^{v/2}$ for a fixed number $v\geq2$ of affine variables, with a matching generic lower bound. Their proof in section 4 uses connected monomial spaces and an even/odd binary-digit split. We do not claim that order as new. The proof below gives the displayed finite integer bound by a different shared-chain construction, records its coefficient-ring scope, and proves the resulting GCT multiplicity exclusion. These are the forms needed here that the cited theorem does not state.
\begin{proof}
Extend $L$ to a basis of linear forms and write it as $x_0$. Dehomogenize $q$ by setting $x_0=1$, obtaining an arbitrary polynomial $p$ of degree at most $m$ in $k-1$ variables. It suffices to construct an affine linear determinant of size $C(m,k)$ equal to $p$. Homogenizing each matrix entry by $x_0$ then gives $x_0^{C(m,k)-m}q$.

For odd $k=2r+1$, split the dehomogenized variables into lists $x$ and $y$ of $r$ variables each. Use the two monomial trees from G-08, each with $\binom{m+r}{r}$ vertices, and connect every pair $(\alpha,\beta)$ with $|\alpha|+|\beta|\leq m$ by its coefficient in $p$. Every source-to-sink path is unique for its two monomial indices, so its path polynomial is $p$. The cofactor construction in G-08 has size one less than the total number of vertices, giving the displayed bound.

For even $k=2r+2$, write the variables as $x_1,\ldots,x_r,y_1,\ldots,y_r,z$ and fix $h$. Decompose the exponent of $z$ uniquely as $e=qh+s$, $0\leq s<h$. On the left, make a directed chain with $Qh$ edges, all labelled $z$, beginning at the source. At its vertex of index $qh$ attach the outward $x$-monomial tree of degrees at most $m-qh$, identifying its constant root with that chain vertex. Every such tree vertex indexed by $(q,\alpha)$ has a unique source path, with product $z^{qh}x^\alpha$. Chain vertices whose indices are not multiples of $h$ are transit vertices and receive no connecting edges. The left graph has
\[
 (Qh+1)+\sum_{q=0}^Q\left[\binom{m-qh+r}{r}-1\right]
 =\sum_{q=0}^Q\binom{m-qh+r}{r}+Q(h-1)
\]
vertices.

On the right, make a chain with vertices $0,\ldots,s_0$, directed toward vertex $0$, the sink, and again label its edges by $z$. At vertex $s$ attach the inward $y$-monomial tree of degrees at most $m-s$, identifying the root with that vertex. Every tree vertex indexed by $(s,\beta)$ has a unique sink path with product $y^\beta z^s$. The right graph has $\sum_{s=0}^{s_0}\binom{m-s+r}{r}$ vertices. Connect left vertex $(q,\alpha)$ to right vertex $(s,\beta)$ with the constant coefficient of $x^\alpha y^\beta z^{qh+s}$ in $p$ when $|\alpha|+|\beta|+qh+s\leq m$.

All connecting edges go from left to right. Thus every source-to-sink path uses exactly one such edge, and the uniqueness of the two tree paths and the division $e=qh+s$ gives exactly the required monomial with its coefficient, once. The path polynomial is $p$. Taking the same cofactor of $I-A$ as in G-08 yields an affine linear determinant of size $C_h(m,k)$. This graph has a path capable of realizing $z^m$ and a connecting edge, so its number of vertices is at least $m+2$; hence $C_h(m,k)\geq m+1$. The odd bound is also greater than $m$. Homogenization therefore uses a nonnegative exponent. Minimizing over the finite choices of $h$ proves the stated representation bound.

For the asymptotic statement in the even case take $h=\lceil\sqrt m\rceil$. The number of binomial summands is at most $m/h+1+h\leq2\sqrt m+2$, each bounded by $\binom{m+r}{r}=O_r(m^r)$, while $Q(h-1)\leq m$. Thus $C(m,k)=O_r(m^{r+1/2})$. The odd case is immediate from its binomial formula. These exponents are $(k-1)/2$ in both cases.

For the GCT consequence restrict a substituted padded form to any fixed $k$ underlying variables. If its restricted padding $L$ is nonzero, apply the determinant construction and pad with $n-C(m,k)$ diagonal copies of $L$. If $L=0$, the form is zero because $n\geq C(m,k)>m$. Thus every restriction lies in $\mathrm{Det}_n$. Apply the highest-weight ideal argument from G-08, which uses exactly this containment, to obtain the multiplicity inequality in every degree.
\end{proof}
\paragraph{Integral scope and growth of the remaining gate.}
With $L=x_0$ the graph and cofactor identity are universal over the coefficient ring. Extending an arbitrary nonzero $L$ to coordinates is a field operation; over a general base one instead works on a chart where a coefficient of $L$ is invertible and chooses the resulting local basis. No global polynomial choice of that basis is asserted. If $n=m^c$ with fixed $c>0$, every fixed $k\geq3$ satisfying $(k-1)/2<c$ is excluded for sufficiently large $m$. This is a quantitative restriction on possible multiplicity obstructions, not a claim that the remaining higher-row obstruction exists.
\begin{proposition}[G-10: An explicit degree bound for the four-row dimension gaps]\label{G-10}
For $n\geq4$ put
\[
 r_n=2n^2+2,\qquad B_n=n^{r_n-1},\qquad
 d_n=r_n(B_n-1)+1.
\]
For every integer $d\geq d_n$, some partition $\lambda\vdash nd$
with exactly four rows satisfies
\[
 m_\lambda(\operatorname{per}_n)>m_\lambda(\det_n).
\]
More generally, the same explicit threshold $d_n$ applies to every
padded pair $(m,n)$ for which the dimension criterion
$\dim\mathcal P_{m,4}+3>r_n$ of R-14 has been established.
The degree bound is very large and does not specify which partition
has the gap. It makes the preceding ``sufficiently large degree''
assertions effective; it gives no larger range of padded matrix sizes.
\end{proposition}
\begin{proof}
Write $Y=\mathbb P(\mathcal D_{n,4})$, $s=\dim Y$, and
$\Delta=\deg Y$. The pencil coefficient map projectivizes to a
rational map
\[
 \phi:\mathbb P^{4n^2-1}\dashrightarrow Y
\]
given by homogeneous polynomials of degree $n$. Its image is dense.
We first check the degree estimate
\begin{equation}\label{G10:image-degree}
 \Delta\leq n^s\leq n^{r_n-1}=B_n.
\end{equation}
Choose $s$ general hyperplanes in the target. Their intersection
with $Y$ consists of $\Delta$ distinct points in a dense open subset
contained in the actual image of $\phi$. Such an open subset exists
because the dominant image is constructible. Each point therefore
has a nonempty fiber in the domain of definition $U$ of $\phi$.
The common zero locus $Z$ of the $s$ pulled-back degree-$n$
equations in $\mathbb P^{4n^2-1}$ meets $U$ in the disjoint union
of these fibers. Every irreducible component of $Z$ meeting $U$
maps to just one of the target points. At least $\Delta$ distinct
components of $Z$ are consequently required. Each has degree at
least one, while repeated projective B\'ezout gives that the sum
of degrees of all components of $Z$ is at most $n^s$.
This proves the first inequality. Components contained in the
base locus only contribute to that upper bound and cannot cover
any of the nonempty fibers in $U$.
The second inequality follows from G-06, which gives $s\leq r_n-1$.
This is the usual image-degree argument with its base locus
accounted for explicitly; compare Sombra (1997), section~1,
printed preprint p.~5, and Proposition~2.10, pp.~11--12,
whose displayed statement treats an everywhere regular map.

The all-degree Hilbert bound of Chardin (1989), the unnumbered
central theorem on printed p.~306, applies to the homogeneous
prime ideal of the reduced irreducible variety $Y$ over
$\mathbb C$. Thus for every $d\geq1$,
\begin{equation}\label{G10:det-Hilbert}
 \dim\mathbb C[\mathcal D_{n,4}]_d
 \leq\Delta\binom{d+s}{s}
 \leq B_n\binom{d+r_n-1}{r_n-1}.
\end{equation}
We use this cited bound rather than assuming that the determinant
coordinate ring is Cohen--Macaulay.

Let $X$ be either $\mathcal P_{n,4}$ or a padded image
$\mathcal Z_{m,n}$ satisfying the stated criterion. It is a
reduced irreducible standard graded affine cone and has dimension
at least $r_n+1$, by R-15 or R-14 respectively. Linear Noether
normalization over the infinite field $\mathbb C$ supplies
$\dim X$ algebraically independent linear coordinate functions.
Any $r_n+1$ of these generate a polynomial subalgebra of its
coordinate ring. Counting its degree-$d$ monomials gives
\begin{equation}\label{G10:per-Hilbert}
 \dim\mathbb C[X]_d\geq\binom{d+r_n}{r_n}.
\end{equation}
The ratio of this binomial to the one in
\eqref{G10:det-Hilbert} is $(d+r_n)/r_n$.
When $d\geq r_n(B_n-1)+1$, that ratio is strictly greater
than $B_n$. Hence the Hilbert function of $X$ is strictly larger
in every such degree.

Decomposing the two degree-$d$ coordinate modules into
$\mathrm{GL}_4$ irreducibles, their dimension difference is
\[
 \sum_{\lambda\vdash nd,\,\ell(\lambda)\leq4}
 \bigl(m_\lambda(X)-m_\lambda(\mathcal D_{n,4})\bigr)
 \dim S_\lambda\mathbb C^4>0.
\]
At least one multiplicity difference is positive. The
highest-weight restriction argument of R-14 identifies these
multiplicities with those of the original orbit closures.
G-05 excludes positive differences with at most three rows.
The remaining partition has exactly four rows, as asserted.
\end{proof}
\paragraph{Effectivity and its limit.}
At the displayed degree, the finite coefficient and isotypic-rank
algorithms of R-05--R-06 can enumerate the partitions and find one
with a positive difference. The proof therefore bounds a finite
search, but its size is not polynomial. In the quartic case the
explicit family at degree $320112+s$ in R-12 is substantially
stronger than this uniform bound.


\begin{proposition}[G-11: Explicit rectangular multiplicity gaps]\label{G-11}
For $n\geq4$ define the following positive integers:
\[
\begin{gathered}
 b_n=\binom{n+3}{3},\qquad \beta_n=6b_n n^{30},\qquad
 T_n=\binom{b_n+\beta_n}{b_n}-1,\\
 E_n=4T_n(\beta_n!),\qquad u_n=2n^2-13,\\
 Q_n=(nE_n)^{u_n-1},\qquad K_n=u_n(Q_n-1)+1.
\end{gathered}
\]
For every integer $k\geq K_n$, put $d=E_n k$ and
$\lambda=(nd/4,nd/4,nd/4,nd/4)$. Then
\begin{equation}\label{G11:rectangle}
 m_\lambda(\operatorname{per}_n)>m_\lambda(\det_n).
\end{equation}
More generally, let $n>m\geq4$. The same conclusion, with
$\operatorname{per}_n$ replaced by
$\ell^{n-m}\operatorname{per}_m$, holds if
\begin{equation}\label{G11:padded-hypotheses}
 \dim\mathcal P_{m,4}+3>2n^2+2
 \quad\hbox{and}\quad n<4m/3.
\end{equation}
In particular it holds for $n=m+1$ and every $m\geq8$, and
for $n=\lfloor\rho m\rfloor$, all sufficiently large $m$,
and any fixed $1<\rho<\sqrt5/2$.
The degree, its progression, and the rectangular partition are
explicit. No useful computational size or polynomial algorithm
is asserted for these bounds.
\end{proposition}
\paragraph{Background used.}
We use the Kempf--Ness closed-orbit criterion in Losev (2006),
Theorem 4, p.~1; the stable-locus and quotient statements of
Mukai (2003), Definition 5.12, Proposition 5.15, Theorem 5.16,
and Corollary 5.17, pp.~165--167; and Derksen (2001),
Theorem 1.1, Proposition 1.2, and Example 2.1, printed p.~956.
The latter gives the displayed bound $\beta_n$ for generators
of the invariant ring of quaternary degree-$n$ forms.
The argument below proves the particular stability witnesses,
Veronese bound, and GCT deduction that those sources do not state.
\begin{proof}
Write $G=\mathrm{SL}_4(\mathbb C)$ and let $V_n$ be the space
of degree-$n$ quaternary forms. First we exhibit stable points
in the cones whose dimensions we compare. The form
\[
 f_n=x_0^n+x_1^n+x_2^n+x_3^n
\]
belongs to both $\mathcal P_{n,4}$ and $\mathcal D_{n,4}$.
Indeed factor $x_0^n+x_1^n$ and $x_2^n+x_3^n$ into binary
linear factors, place them on the diagonal and directed cycle
of an $n\times n$ pencil, and use its two perfect matchings.
For the determinant, multiply one cycle entry by the sign
$(-1)^{n-1}$.

Equip $V_n$ with an $\mathrm{SU}_4$-invariant Hermitian inner
product and use unit monomial vectors. Distinct monomials are
orthogonal. In the normalized sum of the four pure powers,
the mean exponent of each variable is $n/4$. Its moment map
therefore vanishes on the diagonal trace-zero Lie algebra.
It also vanishes on the off-diagonal directions: an operator
$x_j\partial/\partial x_i$, $i\ne j$, cannot carry one of
these pure powers to another when $n\geq3$. Thus the moment
map is zero. The Kempf--Ness criterion makes the $G$-orbit
closed. Its infinitesimal stabilizer is zero, since the
monomials $x_i^{n-1}x_j$ and $x_i^n$ in the derivative of
$f_n$ are distinct for $n\geq3$ and force all matrix entries
to vanish. Its stabilizer is consequently finite. This is a
stable point in each cone.

For the padded case put $a=n-m>0$ and suppose $3a<m$.
Let $e_0$ be the unit monomial $x_0^n$ and $e_i$ the unit
monomial proportional to $x_0^a x_i^m$, $1\leq i\leq3$.
Set
\[
 \theta_0=\frac{m-3a}{4m},\qquad
 \theta_1=\theta_2=\theta_3=\frac{n}{4m},\qquad
 F=\sum_{i=0}^3\sqrt{\theta_i}\,e_i.
\]
All four numbers are positive and sum to one. The mean
exponent of $x_0$ is $n\theta_0+3a\theta_1=n/4$, and
that of $x_i$ is $m\theta_i=n/4$. No two supported monomials
differ by a root $e_i-e_j$, because their differences are
multiples of $m\geq4$. The same moment-map calculation
therefore gives zero. We have $F=x_0^a q$, where
$q=\sum_{i=0}^3 c_i x_i^m$ and all $c_i\ne0$.
The preceding binary-factor construction puts $q$ in
$\mathcal P_{m,4}$, so $F\in\mathcal Z_{m,n}$.

Here is a direct finite-stabilizer check. For a linear vector
field $D_A=\sum_i L_i\partial/\partial x_i$, the equation
$D_A(F)=0$, after division by $x_0^{a-1}$, is
\[
 aL_0q+x_0D_Aq=0.
\]
Reduction modulo $x_0$ shows $L_0=d_0x_0$, because $a\ne0$
and $q(0,x_1,x_2,x_3)\ne0$. It follows that
$D_Aq=-ad_0q$. The distinct mixed monomials
$x_i^{m-1}x_j$ force every off-diagonal entry of $A$ to
vanish. The $x_0^m$ coefficient gives $(m+a)d_0=0$,
and the remaining pure powers give every other diagonal
entry zero. Thus the Lie stabilizer is zero and $F$ is
stable by Kempf--Ness.

For every reduced irreducible cone $X$ just considered,
the stable locus is a nonempty open subset and its quotient
fibers are the $15$-dimensional $G$-orbits (Mukai, loc. cit.).
Hence the dimension-of-fibers theorem yields
\begin{equation}\label{G11:quotient-dimension}
 \dim\mathbb C[X]^G=\dim X-15.
\end{equation}
In particular, writing $A_D=\mathbb C[\mathcal D_{n,4}]^G$
and $A_X=\mathbb C[X]^G$, G-06, R-15, and the first
hypothesis in \eqref{G11:padded-hypotheses} imply
\[
 \dim A_D\leq u_n,\qquad \dim A_X\geq u_n+1.
\]
The invariant functor is exact in characteristic zero.
Thus both invariant rings are graded quotients of
$\mathbb C[V_n]^G$, including when $X$ is padded; see also
Mukai, Proposition 5.10 and its proof, pp.~164--165.
Derksen's bound, specialized to $q=4$ and symmetric degree
$n$, says that this ambient invariant ring is generated in
degrees at most $\beta_n=6b_n n^{30}$. The dual convention
for forms makes no difference: the contragredient
transforms the action by an automorphism of $\mathrm{SL}_4$.
We may take at most $T_n$ homogeneous generators, because
the vector space of all coefficient polynomials in $b_n$
variables of positive degree at most $\beta_n$ has dimension
$T_n$.

We check explicitly that the $E_n$th Veronese of each ring
is generated in its first degree. More generally, suppose
an algebra has $r\leq T_n$ homogeneous generators with
positive degrees $d_i\leq\beta_n$, and put
$B_0=\beta_n!$. In a monomial in these generators of total
degree $kE_n$, remove as many pure blocks of
$B_0/d_i$ copies of generator $i$ as possible. Every removed
block has degree $B_0$. The remaining exponents are less
than $B_0/d_i$, so the residual degree is less than
$rB_0\leq T_n B_0$. It is an integer multiple $hB_0$,
because the original and removed degrees are such multiples;
thus $0\leq h<T_n$. There are $4kT_n-h$ pure blocks.
Combine the residual monomial with $4T_n-h$ of these blocks
to form a monomial of degree $E_n=4T_nB_0$, then divide
the other $4(k-1)T_n$ blocks into groups of $4T_n$.
This factors the original monomial into $k$ monomials of
degree $E_n$. Applying this to spanning monomials proves
the claim, also after passing to a quotient.

Consequently $A_D^{(E_n)}$ and $A_X^{(E_n)}$ are standard
graded domains, with dimensions unchanged. Set
$s=\dim\operatorname{Proj}A_D^{(E_n)}\leq u_n-1$.
This projective dimension is positive: block padding embeds
$\mathcal D_{4,4}$, whose dimension is $34$ by G-07,
into $\mathcal D_{n,4}$. By \eqref{G11:quotient-dimension}
the latter invariant ring has dimension at least $19$,
so $s\geq18$.
Choose a basis of its degree-one piece represented by
homogeneous invariant polynomials of degree $E_n$ on $V_n$.
Their composition with the determinant-pencil map is a
collection of homogeneous polynomials of degree $nE_n$.
They give a dominant rational map from
$\mathbb P^{4n^2-1}$ to $\operatorname{Proj}A_D^{(E_n)}$.
Dominance follows from dominance of the pencil map and
surjectivity of the affine quotient (Mukai, Theorem 5.9);
the Veronese morphism has this projective image. G-10's
base-locus-aware image-degree argument gives
\[
 \deg\operatorname{Proj}A_D^{(E_n)}
 \leq(nE_n)^s\leq Q_n.
\]
Chardin's bound in G-10 therefore gives, for every $k\geq1$,
\[
 \dim(A_D)_{E_nk}
 \leq Q_n\binom{k+u_n-1}{u_n-1}.
\]
Linear Noether normalization of the standard graded ring
$A_X^{(E_n)}$, whose dimension is at least $u_n+1$, gives
\[
 \dim(A_X)_{E_nk}\geq\binom{k+u_n}{u_n}.
\]
The ratio of the binomials is $(k+u_n)/u_n>Q_n$ when
$k\geq K_n$. Thus the first invariant dimension is strictly
smaller than the second.

Finally, in coefficient degree $d$, every polynomial
$\mathrm{GL}_4$ weight has size $nd$. The only irreducible
polynomial representation with trivial restriction to
$\mathrm{SL}_4$ is a determinant power; it has shape
$(nd/4)^4$, when $4\mid nd$, and is one-dimensional.
Here $E_n$ is divisible by four. The invariant dimension
difference is therefore exactly the difference of the two
multiplicities for the displayed rectangular partition.
The highest-weight restriction argument of G-05 and R-14
identifies these with the full orbit-closure multiplicities.
This proves \eqref{G11:rectangle} and its padded version.
R-14 and R-16 supply the stated particular ranges; all
satisfy $n<4m/3$.
\end{proof}
\paragraph{Relation to the rectangular bridge.}
The Haar projector onto $\mathrm{SU}_4$ invariants isolates
precisely these rectangular highest weights. Equivalently,
the corresponding character integral is the dimension of
this invariant graded piece. The proof makes that positive
rank exceed the determinant rank; it does not replace an
orbit-closure multiplicity by an ambient Kronecker
coefficient. The quartic formula G-07/R-12 is much sharper
where it applies. These deductions make no literature-priority
claim and are independent of RH and either proposed GRH input.

\begin{proposition}[G-12: A padding exclusion for the rectangular sector]\label{G-12}
Let $m\geq1$, $n>m$, and $n\geq4m/3$. For every degree-$m$
form $f$, every coordinate-ring degree $d\geq1$ with
$4\mid nd$, and $\lambda=(nd/4)^4$, one has
\begin{equation}\label{G12:exclusion}
 m_\lambda(\ell^{n-m}f)\leq m_\lambda(\det_n).
\end{equation}
If $n>4m/3$, the multiplicity on the left is zero.
The equality case $n=4m/3$ is included in the inequality.
Thus the rectangular gaps of G-11 cannot extend to or beyond
this padding boundary, even though nonrectangular four-row
questions remain.
\end{proposition}
\begin{proof}
As in G-05, restrict substituted forms to four variables.
They have shape $L^a q$, where $a=n-m$ and $q$ has degree
$m$. If $L=0$, the form is zero. Otherwise a change of
coordinates takes $L$ to $x_0$; it suffices to use a
$\mathrm{GL}_4$ change, since the determinant cone and each
homogeneous invariant ideal are stable under this group.
Act on the variables by the one-parameter subgroup
\[
 \eta(t)=\operatorname{diag}(t^3,t^{-1},t^{-1},t^{-1})
 \quad\hbox{of }\mathrm{SL}_4.
\]
A monomial in $x_0^a q$ with $x_0$ exponent $a+j$ has
weight $3a-m+4j$. If $3a>m$, all these weights are positive,
so $\eta(t)x_0^a q\longrightarrow0$ as $t\longrightarrow0$.
Every homogeneous positive-degree invariant therefore
vanishes on every such padded form and on their orbit
closure. Exactness of invariants identifies its invariant
coordinate ring with $\mathbb C$, whose positive-degree pieces vanish,
which proves the zero assertion.

If $3a=m$, the same limit instead equals
\[
 F_0=x_0^a q(0,x_1,x_2,x_3).
\]
Every ternary degree-$m$ form belongs to $\mathcal D_{m,3}$
by G-05 (including the elementary degrees one and two).
Embedding it and adding $a$ diagonal copies of $x_0$
shows $F_0\in\mathcal D_{n,4}$; the argument also holds for
limits because the block-padding map is polynomial.
Let $P$ be any homogeneous $\mathrm{SL}_4$ invariant
vanishing on $\mathcal D_{n,4}$. Since $P$ is constant
along $\eta$, continuity gives
$P(x_0^a q)=P(F_0)=0$. Thus the determinant ideal intersected
with the ambient invariant ring is contained in the ideal
of the restricted padded orbit closure intersected with
that ring. Both invariant coordinate rings are quotients
of the same ambient invariant ring, by exactness of
invariants. This containment gives a surjection from the
determinant invariant coordinate ring to the padded one,
and hence the degreewise dimension inequality.
Invariant dimensions are precisely the rectangular
multiplicities, as in G-11; highest-weight restriction
transfers the inequality to the full orbit closures.
\end{proof}
\paragraph{Source and scope.}
The strict inequality is also the rectangular specialization
of the usual padding weight constraint, but the proof above
gives its one-parameter-subgroup form and handles the
boundary by the ternary determinant theorem. This is an
exclusion for this specified rectangular sector. It does
not compare all four-row multiplicities or arbitrary
Kronecker coefficients.

\begin{proposition}[G-13: An explicit finite bound for the first rectangular gap]\label{G-13}
For $n\geq4$ put
\[
 \beta=6\binom{n+3}{3}n^{30},\qquad u=2n^2-13,\qquad
 C=(n\beta^2)^u,
\]
and define
\begin{equation}\label{G13:first-bound}
 s_*=(u+1)(C-1)+1,\qquad d_* = \beta s_*.
\end{equation}
There is an integer $d$ with $1\leq d\leq d_*$ and $4\mid nd$
such that, for $\lambda=(nd/4)^4$,
\[
 m_\lambda(\operatorname{per}_n)>m_\lambda(\det_n).
\]
The same conclusion holds with $\operatorname{per}_n$ replaced by
$\ell^{n-m}\operatorname{per}_m$ whenever $n>m\geq4$ and
\[
 \dim\mathcal P_{m,4}+3>2n^2+2,
 \qquad n<4m/3.
\]
In particular it applies to $n=m+1$, $m\geq8$, and, by R-18,
to $n=\lfloor\rho m\rfloor$ for all sufficiently large $m$
and each fixed $1<\rho<\sqrt{3/2}$.

The bound satisfies $\log d_*=O(n^2\log n)$. This is an
existence bound within a finite interval. It does not specify
the individual successful degree or replace G-11's assertion
at every degree of its displayed progression.
\end{proposition}

\paragraph{Background and scope.}
We use the stable points, exact invariant quotients, and quotient
dimensions established in G-11, together with Derksen (2001),
Theorem~1.1 and Example~2.1, printed p.~956. We also use the
projective image-degree argument of G-10 and Chardin's
unnumbered main theorem (1989), printed p.~306. The argument
below compares a weighted degree filtration with the ordinary
degree filtration in an affine generator embedding; it needs
no factorial Veronese.

\begin{proof}
Let $G=\mathrm{SL}_4(\mathbb C)$ and let $V_n$ be the space
of quaternary forms of degree $n$. Write
\[
 A_D=\mathbb C[\mathcal D_{n,4}]^G,
 \qquad A_X=\mathbb C[X]^G,
\]
where $X$ is $\mathcal P_{n,4}$ or the padded cone in the
statement. The stable-point and dimension arguments of G-11
give
\begin{equation}\label{G13:dims}
 t:=\dim A_D\leq u,\qquad \dim A_X\geq u+1.
\end{equation}
These are positively graded domains, with degree-zero part
$\mathbb C$, and both are graded quotients of
$\mathbb C[V_n]^G$. Choose homogeneous generators
$g_1,\ldots,g_r$ of the latter ring with positive degrees
$e_i\leq\beta$. Their images generate both quotient rings.

First bound the cumulative graded dimension of $A_D$. The
images of the generators give a closed affine embedding
\[
 Y_D=\operatorname{Spec} A_D\lhook\joinrel\longrightarrow
 \mathbb A^r.
\]
Compose this embedding with the determinant-pencil map and
the affine quotient. The resulting map from
$\mathbb A^{4n^2}$ has coordinates
$g_i(\det L)$, homogeneous of degrees $ne_i\leq n\beta$
in the pencil parameters. Its image is dense in $Y_D$.
To homogenize it, adjoin a source coordinate $z$ and use
the degree-$n\beta$ coordinate polynomials
\[
 z^{n\beta},\qquad
 z^{n\beta-ne_i}g_i(\det L),\quad 1\leq i\leq r.
\]
They define a rational map from $\mathbb P^{4n^2}$ whose
image closure is the ordinary projective closure
$\overline{Y_D}\subseteq\mathbb P^r$. It has dimension
$t$. G-10's image-degree bound, allowing a base locus, gives
\begin{equation}\label{G13:affine-degree}
 \deg\overline{Y_D}\leq(n\beta)^t\leq(n\beta)^u.
\end{equation}

For $d\geq0$, let
\[
 H_D^{\leq}(d)=\sum_{j=0}^d\dim(A_D)_j.
\]
Every element of $(A_D)_j$ has a preimage homogeneous of
weighted degree $j$ in $\mathbb C[T_1,\ldots,T_r]$ with
$\deg T_i=e_i$. Since every $e_i\geq1$, each monomial in
that preimage has ordinary degree at most $j$. It follows
that the direct sum of the graded pieces through $d$ is
contained in the space of functions on $Y_D$ represented
by ordinary polynomials of degree at most $d$.
This is an inclusion of filtered vector spaces, not an
identification of the two gradings.

The latter ordinary filtered dimension equals the
degree-$d$ Hilbert function of $\overline{Y_D}$:
homogenization with its first coordinate, followed by
dehomogenization on the dense affine chart, identifies the
two quotients. The projective closure is integral over
$\mathbb C$. Also $t\geq19$ by the fixed block-padding and
stable-quotient observation in G-11, so the positive-dimension
hypothesis in Chardin's stated theorem holds. That theorem
and \eqref{G13:affine-degree} yield
\begin{equation}\label{G13:upper}
 H_D^{\leq}(d)
 \leq(n\beta)^t\binom{d+t}{t}
 \leq(n\beta)^u\binom{d+u}{u}.
\end{equation}

For the other ring, a maximal algebraically independent
subset of the images of $g_1,\ldots,g_r$ has size
$\dim A_X$. Indeed those images generate its fraction
field, and maximality makes every remaining generator
algebraic over the chosen subset. By \eqref{G13:dims},
choose $u+1$ algebraically independent images. Each has
degree at most $\beta$. Their monomials of ordinary total
degree at most $s$ are linearly independent and have
graded degree at most $\beta s$. Consequently
\begin{equation}\label{G13:lower}
 H_X^{\leq}(\beta s)
 :=\sum_{j=0}^{\beta s}\dim(A_X)_j
 \geq\binom{s+u+1}{u+1}.
\end{equation}

For every nonnegative integer $s$, the exact binomial ratio
satisfies
\begin{align*}
 \frac{\binom{s+u+1}{u+1}}{\binom{\beta s+u}{u}}
 &=\frac{s+u+1}{u+1}
     \prod_{i=1}^u\frac{s+i}{\beta s+i}\\
 &\geq\frac{s+u+1}{(u+1)\beta^u}.
\end{align*}
The last inequality follows term by term from
$\beta(s+i)\geq\beta s+i$. At $s=s_*$ one has
\[
 s_*+u+1=(u+1)C+1,
\]
so the final lower bound is strictly greater than
$C/\beta^u=(n\beta)^u$. Equations
\eqref{G13:upper}--\eqref{G13:lower} therefore give
\[
 H_X^{\leq}(d_*)>H_D^{\leq}(d_*).
\]
The degree-zero dimensions are both one. Hence for some
$1\leq d\leq d_*$ the individual graded pieces satisfy
$\dim(A_X)_d>\dim(A_D)_d$.

As in G-11, the only polynomial $\mathrm{GL}_4$
representation of weight size $nd$ that is trivial on
$\mathrm{SL}_4$ is the one-dimensional determinant power
of shape $(nd/4)^4$. If $4\nmid nd$, both invariant
pieces are zero. Thus the successful degree necessarily
satisfies $4\mid nd$, and the invariant dimension gap is
the desired rectangular multiplicity gap. The first-four-
variable highest-weight argument in G-05/R-14 transfers
it to the full orbit closures.

Finally $\beta=O(n^{33})$ and $u=O(n^2)$, so the explicit
formula \eqref{G13:first-bound} gives
$\log d_*=O(n^2\log n)$. The particular padded ranges
follow from R-14/R-16 and R-18, respectively; their
dimension gaps and $n<4m/3$ are the same hypotheses already
used to construct the stable padded point in G-11.
\end{proof}

\paragraph{Complexity lower-bound context.}
The multiplicity statements G-10--G-13 locate particular representation
sectors and degrees. They do not improve the established quadratic
lower bound $\overline{\operatorname{dc}}(\operatorname{per}_m)
\geq m^2/2$ of Landsberg, Manivel, and Ressayre (2013);
the inspected April 2010 author version states it as Theorem~1.0.1,
p.~1 (G-B16). No claim about the strongest currently proposed lower
bound is needed here. A noncontainment result by itself does not
specify the four-row or rectangular multiplicity differences proved
in the present restricted parameter ranges.

\begin{proposition}[G-14: First-jet dominance and a stronger padding boundary]\label{G-14}
Let $m\geq1$, and write a quaternary form as
\[
 q(x_0,x_1,x_2,x_3)=\sum_{j=0}^m x_0^j q_j(x_1,x_2,x_3).
\]
The first-jet map
\[
 j_1:\mathcal D_{m,4}\longrightarrow
   \operatorname{Sym}^m\mathbb C^3\oplus
   \operatorname{Sym}^{m-1}\mathbb C^3,
 \qquad q\longmapsto(q_0,q_1),
\]
has Zariski-dense image.

Consequently let $n>m\geq1$, put $a=n-m$, and let $f$ be
any degree-$m$ form in the common ambient GCT setting.
For every $d\geq1$ and every partition $\lambda\vdash nd$
with at most four rows and
\begin{equation}\label{G14:boundary}
 \lambda_1\leq ad+1,
\end{equation}
one has
\[
 m_\lambda(\ell^a f)\leq m_\lambda(\det_n).
\]
If $\lambda_1<ad$, the multiplicity on the left is zero.
Thus every positive multiplicity gap in this sector must
satisfy $\lambda_1\geq ad+2$.
\end{proposition}

\paragraph{Scope and background.}
The proof uses the coefficient grading and first-four-variable
highest-weight restriction of G-05, together with elementary
cyclic cofactors and the finite geometric-product identity
already proved in R-16. It requires characteristic zero.
The jet-dominance calculation is given in full; no higher-jet
dominance or Lefschetz assertion is assumed.

\begin{proof}
The case $m=1$ of the first assertion is immediate. Suppose
$m\geq2$, choose a primitive $m$th root of unity $\zeta$,
and let $A$ be the ternary cyclic pencil with diagonal
entries
\[
 \ell_i=x_1-\zeta^i x_2,\qquad 0\leq i<m,
\]
and entries $x_3$ in positions $(i,i+1)$, with indices
modulo $m$. Its determinant has the two perfect matchings
and equals
\[
 x_1^m-x_2^m+(-1)^{m-1}x_3^m.
\]
We claim that the cofactors of $A$ span every ternary form
of degree $m-1$.

For each $0\leq k\leq m-1$, the cofactors using exactly
$k$ cycle entries are, up to nonzero signs, $x_3^k$ times
the $m$ cyclic translates of a consecutive product of
$h=m-1-k$ of the binary forms $\ell_i$. This follows from
the unique matching after deleting the relevant row and
column. A consecutive product, after a cyclic shift,
has the form
\[
 F_h(x_1,x_2)=\prod_{i=0}^{h-1}(x_1-\zeta^i x_2).
\]
Every coefficient in it is nonzero. Indeed its coefficient
of $x_1^{h-r}x_2^r$ is
\[
 (-1)^r\zeta^{r(r-1)/2}
   \prod_{j=1}^r\frac{1-\zeta^{h-r+j}}{1-\zeta^j},
 \qquad 0\leq r\leq h,
\]
with the empty product interpreted as one. Every exponent
in a nonempty numerator or denominator lies between $1$
and $h<m$. Thus no factor is zero. The cyclic substitution
$x_2\mapsto\zeta x_2$ has distinct characters on the
monomials with $0\leq r\leq h$. Fourier projection of the
$m$ translates therefore spans every binary form of degree
$h$. Varying $k$ proves the claim, since
\[
 \operatorname{Sym}^{m-1}\mathbb C^3
 =\bigoplus_{k=0}^{m-1}x_3^k
        \operatorname{Sym}^{m-1-k}\langle x_1,x_2\rangle.
\]

Now consider the parameter map
\[
 (A',B)\longmapsto j_1\det(A'+x_0B),
\]
where $A'$ is an arbitrary ternary linear pencil and $B$
is a constant $m\times m$ matrix. At the point $(A,0)$,
the derivative in the $A'$ directions changes $q_0$ by
arbitrary linear-form multiples of the cofactors of $A$.
By the claim these span all of
$\operatorname{Sym}^m\mathbb C^3$. It does not change
$q_1$, since the coefficient of $x_0$ is linear in $B$.
The derivative in the $B$ directions changes $q_1$ by
arbitrary constant linear combinations of the cofactors,
which span $\operatorname{Sym}^{m-1}\mathbb C^3$, and
does not change $q_0$. The differential is consequently
surjective onto the direct sum. The characteristic-zero
Jacobian criterion for dominance proves the first assertion.

For the multiplicity consequence, work first in four
variables and use the positive coefficient weight convention
of G-05/G-11 (equivalently, dualize the coordinate-ring
representations consistently). In a polynomial irreducible
of shape $\lambda$, every weight has each coordinate at
most $\lambda_1$. For example, the Schur construction uses
$\lambda_1$ exterior-power columns, and a given basis vector
can occur at most once in each column. This proves the
required bound on the $x_0$ coordinate of every weight.

Let $R=\mathbb C[\operatorname{Sym}^n\mathbb C^4]$ be
the coefficient polynomial ring, and let $R_{d,\lambda}$
be its degree-$d$ $\lambda$-isotypic subspace. Take
\[
 P\in I(\mathcal D_{n,4})\cap R_{d,\lambda}.
\]
The weight bound says that every coefficient monomial
occurring in $P$ has total $x_0$ weight at most
$\lambda_1$. Evaluate $P$ at a form $x_0^a q$. A
coefficient coming from $q_j$ has $x_0$ weight $a+j$;
a product of $d$ coefficients therefore has weight
\[
 ad+j_1+\cdots+j_d.
\]
All coefficient variables of smaller individual $x_0$
weight vanish on $x_0^a q$. If $\lambda_1<ad$, nothing
survives and every polynomial in $R_{d,\lambda}$ vanishes
on every such padded form.

Under \eqref{G14:boundary}, the restriction
$P(x_0^a q)$ depends only on $q_0,q_1$, because every
surviving coefficient monomial has
$j_1+\cdots+j_d\leq1$. It is a polynomial function of
these first-jet data. For every $q\in\mathcal D_{m,4}$,
block padding gives $x_0^a q\in\mathcal D_{n,4}$,
including the closure points. Hence this jet polynomial
vanishes on the dense image of $j_1$ just proved, and
therefore vanishes for every pair $(q_0,q_1)$. It follows
that $P(x_0^a q)=0$ for every quaternary form $q$ of
degree $m$.

To allow any padding form $L$, use a $\mathrm{GL}_4$
change taking $L$ to $x_0$. The determinant ideal and
the isotypic subspace are both stable under that group,
so the preceding conclusion applies to the corresponding
translate of $P$. Thus $P$ vanishes on every $L^a q$,
and in particular on the restricted padded orbit closure
associated with $\ell^a f$.
We have proved the containment of the determinant-ideal
$\lambda$-isotypic part in the corresponding padded ideal
part. Passing to the characteristic-zero semisimple
degree-$d$ quotients gives the stated multiplicity
inequality. The first-four-variable highest-weight
restriction identifies these multiplicities with those
in the full ambient orbit closures, as in G-05/R-14.
The same argument gives the zero assertion when
$\lambda_1<ad$.
\end{proof}

\paragraph{A precisely stated higher-jet continuation.}
For $0\leq h\leq m$, suppose that the projection
\[
 j_h:\mathcal D_{m,4}\longrightarrow
   \bigoplus_{j=0}^h\operatorname{Sym}^{m-j}\mathbb C^3
\]
has dense image. The identical coefficient-weight argument
then proves the same multiplicity inequality whenever
$\lambda_1\leq ad+h$: every surviving coefficient
monomial depends only on those jets. This is a conditional
continuation with an explicit geometric hypothesis.
The first-jet case above is proved without that hypothesis.
No dominance assertion for $h=2$ or $h=3$ is made here.

\begin{proposition}[G-15: Uniform third-jet dominance]\label{G-15}
For every $m\geq1$, the map
\[
 j_{\leq3}:\mathcal D_{m,4}\longrightarrow
   \bigoplus_{j=0}^{\min\{3,m\}}
          \operatorname{Sym}^{m-j}\mathbb C^3,
 \qquad \sum_{j=0}^m t^j q_j\longmapsto(q_0,q_1,q_2,q_3),
\]
with nonexistent components omitted, has Zariski-dense image.
For $m\geq4$, order three is the largest order for which the
full jet map can be dominant.

Consequently, if $n>m\geq1$, $a=n-m$, $d\geq1$, and
$\lambda\vdash nd$ has at most four rows, then
\[
 \lambda_1\leq ad+3
 \quad\Longrightarrow\quad
 m_\lambda(\ell^a f)\leq m_\lambda(\det_n)
\]
for every degree-$m$ form $f$ in the common ambient setting.
Every positive gap in this sector therefore satisfies
$\lambda_1\geq ad+4$. In particular a rectangular gap
with $\lambda=(nd/4)^4$ requires
\[
 d(4m-3n)>12.
\]
Since $4\mid nd$ for this rectangular shape, the last
quantity is divisible by four; equivalently,
$d(4m-3n)\geq16$ is necessary.
\end{proposition}

\paragraph{Background and proof strategy.}
We use the standard $\mathfrak{sl}_2$ decompositions
\[
 \operatorname{End}(\operatorname{Sym}^{m-1}\mathbb C^2)
       =\bigoplus_{l=0}^{m-1}\operatorname{Sym}^{2l}\mathbb C^2,
 \qquad
 \operatorname{Sym}^{r}(\operatorname{Sym}^2\mathbb C^2)
       =\bigoplus_{b=0}^{\lfloor r/2\rfloor}
             \operatorname{Sym}^{2r-4b}\mathbb C^2.
\]
These are the Clebsch--Gordan and harmonic decompositions;
see Fulton and Harris (1991), section~11.2, Exercises~11.11
and 11.14, printed pp.~151 and 153. We specify their
covariants, the finite orthogonal-polynomial calculation,
and the required rank argument below. No generic strong
Lefschetz property or Gulliksen--Neg\aa rd theorem is
assumed in this proof.

\begin{proof}
The case $m=1$ is immediate. Fix $m\geq2$ and let
$\rho$ be the irreducible representation of
$\mathfrak{sl}_2$ on $V=\operatorname{Sym}^{m-1}\mathbb C^2$.
Use standard generators $H,E,F$, with
$[H,E]=2E$, $[H,F]=-2F$, and $[E,F]=H$.
Choose the invariant Hermitian metric for which
$\rho(H)$ is Hermitian and $\rho(E)^*=\rho(F)$.
Its $m$ distinct $H$-eigenvalues are
\[
 \lambda_j=m-1-2j,\qquad 0\leq j<m.
\]
For $z=(z_1,z_2,z_3)$ put
\[
 A(z)=z_1H+z_2(E+F)+\mathrm i z_3(E-F),\qquad
 q=z_1^2+z_2^2+z_3^2,
\]
and consider the explicit linear pencil
\begin{equation}\label{G15:pencil}
 M(t,z)=tI_m-\rho(A(z)).
\end{equation}
The invariant quadratic form is normalized so that
$q(H)=1$. In the defining two-dimensional representation,
$A(z)^2=q(z)I_2$.

Let $\mathcal H_l$ be the space of harmonic homogeneous
polynomials of degree $l$ in $z_1,z_2,z_3$. It is the
irreducible representation of highest weight $2l$ and
dimension $2l+1$. The harmonic decomposition reads
\begin{equation}\label{G15:harmonic}
 \mathbb C[z_1,z_2,z_3]
     =\bigoplus_{l\geq0}\mathcal H_l\otimes\mathbb C[q].
\end{equation}
One can obtain this directly by applying the Laplacian to
$q^b h_l$: for harmonic $h_l$,
\[
 \Delta(q^b h_l)=2b(2l+2b+1)q^{b-1}h_l.
\]
This gives the direct decomposition and its dimensions in
characteristic zero. The indicated irreducibility is the
standard highest-weight calculation for the adjoint
three-dimensional representation.

\paragraph{1. Canonical harmonic matrix covariants.}
Write the multiplicity-free decomposition of
$\operatorname{End}(V)$ as
$\bigoplus_{l=0}^{m-1}\mathcal E_l$, where
$\mathcal E_l\simeq\mathcal H_l$, and let $\pi_l$ be its
equivariant projection. The trace Hermitian form makes
these distinct summands orthogonal. The polynomial map
\[
 A\longmapsto\pi_l\bigl(\rho(A)^l\bigr)
\]
is nonzero and has harmonic coefficient functions of
degree $l$. Indeed the symmetrized multiplication map
from the $l$th symmetric power of the adjoint
representation to $\operatorname{End}(V)$ is equivariant.
Its highest-weight vector $E^l$ maps to
$\rho(E)^l\ne0$ for $l<m$, and this vector has highest
weight $2l$, so lies in $\mathcal E_l$. The harmonic
decomposition says that this is the unique copy of
highest weight $2l$ in symmetric degree $l$.

The value of this covariant at $H$ is nonzero. Otherwise
equivariance and homogeneity would make it vanish on
every scalar multiple of the semisimple orbit of $H$,
which is the dense open set $q\ne0$, contradicting its
nonvanishing. This value is diagonal, since it has
Cartan weight zero.

These diagonal vectors have an elementary description.
Put the uniform positive inner product on functions on
$\{\lambda_0,\ldots,\lambda_{m-1}\}$:
\[
 \langle f,g\rangle=\frac1m\sum_{j=0}^{m-1}
                          f(\lambda_j)\overline{g(\lambda_j)}.
\]
Let $p_0,\ldots,p_{m-1}$ be the real orthonormal
polynomials obtained from $1,x,\ldots,x^{m-1}$, with
positive leading coefficients; thus $p_0=1$.
The zero-weight vector in $\mathcal E_l$ is proportional
to $p_l(\rho(H))$. To check the assertion, note that
$\rho(H)^r$ lies in the sum of the zero-weight spaces of
$\mathcal E_0,\ldots,\mathcal E_r$ by the symmetric-power
decomposition. The top projection is nonzero by the
preceding argument. Induction therefore identifies their
successive spans with those of $1,x,\ldots,x^r$ on the
spectrum. Orthogonality of the endomorphism summands
is precisely orthogonality for the displayed uniform
inner product. This proves the assertion, including
its degree and nonzero leading coefficient.

Normalize the covariant $T_l(A)$ so that
$T_l(H)=p_l(\rho(H))$. Its coefficient functions span
all of $\mathcal H_l$, since it is a nonzero map between
the corresponding irreducible representations.
Equivariance and \eqref{G15:harmonic} now give uniquely
\begin{equation}\label{G15:adjexpansion}
 \pi_l\operatorname{adj} M(t,z)=T_l(A(z))F_l(t,q),
 \qquad 0\leq l<m,
\end{equation}
where $F_l$ is a polynomial homogeneous of weighted
degree $m-1-l$ for $\deg t=1$, $\deg q=2$.

\paragraph{2. The scalar factors are trailing Jacobi determinants.}
Multiplication by $x$ on the functions on the $m$ nodes,
expressed in the orthonormal basis $p_0,\ldots,p_{m-1}$,
is a real symmetric tridiagonal matrix $J$.
Tridiagonality follows because $xp_l$ is orthogonal
to every $p_r$ with $r<l-1$. Its off-diagonal entries
$b_1,\ldots,b_{m-1}$ are strictly positive: they are
the ratios of the positive leading coefficients of
successive orthonormal polynomials. The nodes and
measure are symmetric under $x\mapsto-x$, so $p_l$
has parity $l$. Thus every diagonal entry
$\langle xp_l,p_l\rangle$ is zero.

Let $J^{[l+1]}$ denote the trailing principal submatrix
on indices $l+1,\ldots,m-1$, with the empty determinant
equal to one, and define
\[
 Q_l(t)=\det(tI-J^{[l+1]}),\qquad 0\leq l<m.
\]
The spectrum of $J$ is the node set. If
$D(t)=\prod_j(t-\lambda_j)$, evaluating
\eqref{G15:adjexpansion} at $H$ and taking its
orthonormal trace coefficient gives
\[
 F_l(t,1)=\frac1m\sum_jp_l(\lambda_j)
                         \frac{D(t)}{t-\lambda_j}.
\]
On the other hand the spectral expression for the
$(l,0)$ entry of $(tI-J)^{-1}$ is
\[
 \frac1m\sum_j\frac{p_l(\lambda_j)}{t-\lambda_j}.
\]
Expansion of the corresponding tridiagonal cofactor
shows that its numerator is
$b_1\cdots b_l Q_l(t)$: the unique path from index
zero to index $l$ contributes the indicated product,
and all remaining rows and columns give the trailing
principal determinant. Therefore
\begin{equation}\label{G15:jacobi}
 F_l(t,1)=b_1\cdots b_l Q_l(t).
\end{equation}
In particular every factor is nonzero. Weighted
homogeneity then determines $F_l(t,q)$ from this
univariate identity. All $Q_l$ have the parity of
their degree $m-1-l$.

We record the elementary spectral facts needed below.
Every nonempty trailing block is real symmetric,
tridiagonal, and has positive off-diagonal entries.
An eigenvector is determined by its first coordinate
using the three-term recurrence; that coordinate cannot
be zero. Hence all eigenvalues are simple. Conjugation
by the alternating-sign diagonal matrix negates the
block, so its eigenvalues occur in opposite pairs.
An even-sized block has no zero eigenvalue, as its
determinant at zero is a nonzero signed product of
alternate squared off-diagonal entries. An odd-sized
block has a simple zero eigenvalue. These assertions
also hold for blocks of sizes zero or one with the
obvious interpretations.

\paragraph{3. Cofactor multiples span every third jet.}
Let $S=\mathbb C[t,z_1,z_2,z_3]$, and let $\mathcal C_m$
be the degree-$m$ space spanned by all linear-form
multiples of the entries of $\operatorname{adj} M$.
It is the differential image of the determinant-pencil
map at \eqref{G15:pencil}. Formula
\eqref{G15:adjexpansion} implies that it contains
\[
 tF_l(t,q)\mathcal H_l\qquad(0\leq l<m).
\]
It also contains
\begin{equation}\label{G15:lowerproduct}
 qF_{l+1}(t,q)\mathcal H_l\qquad(0\leq l\leq m-2).
\end{equation}
For clarity, multiplication of a harmonic polynomial of
degree $l+1$ by a linear form has a nonzero projection
onto $q\mathcal H_l$. Indeed
$\Delta(z_i h)=2\partial_i h$; these derivatives span
$\mathcal H_l$ by equivariance and nonvanishing.
The space $\mathcal C_m$ is invariant, so projecting
its products onto this irreducible summand stays in
$\mathcal C_m$. Multiplication by the invariant
$F_{l+1}$ gives \eqref{G15:lowerproduct}.
Similarly the top harmonic component of
$\mathcal H_1\mathcal H_{m-1}$ supplies
$\mathcal H_m\subseteq\mathcal C_m$, since
$F_{m-1}$ is a nonzero constant.

The quotient $S_m/t^4S_{m-4}$ is exactly the third-jet
space, with $S_{m-4}=0$ if $m<4$. Its harmonic
decomposition has, for each $0\leq l\leq m$, the
multiplicity space
\begin{equation}\label{G15:targetspace}
 \operatorname{span}\{t^j q^{(D-j)/2}:
         0\leq j\leq\min\{3,D\},\ j\equiv D\pmod2\},
 \qquad D=m-l,
\end{equation}
tensored with $\mathcal H_l$. It has dimension one
for $D=0,1$, and dimension two for $D\geq2$.
The cases $D=0$ and $D=1$ are already supplied by
$\mathcal H_m$ and $tF_{m-1}\mathcal H_{m-1}$.

For $D\geq2$, specialize $q=1$. On the fixed
multiplicity space \eqref{G15:targetspace}, this is
an isomorphism onto the corresponding powers of $t$
modulo $t^4$. By \eqref{G15:jacobi}, it suffices to
show that $tQ_l(t)$ and $Q_{l+1}(t)$ are independent
modulo $t^4$.

If $D$ is even, $Q_l$ has odd degree $D-1$ and a
simple root at zero, while $Q_{l+1}$ has even degree
$D-2$ and nonzero value at zero. Thus $tQ_l$ starts
with a nonzero $t^2$ term, while $Q_{l+1}$ has a
nonzero constant term. They are independent in the
even two-dimensional space modulo $t^4$.

If $D$ is odd, then $D\geq3$. The block
$K=J^{[l+1]}$ is real symmetric and invertible,
of even size $D-1$. Write $e$ for its first
coordinate vector. The diagonal resolvent identity gives
\[
 \frac{Q_{l+1}(t)}{Q_l(t)}
    =e^{\mathsf T}(tI-K)^{-1}e.
\]
The alternating-sign symmetry fixes $e$ and negates
$K$, so all diagonal entries of odd powers of
$K^{-1}$ vanish. Expanding the inverse at zero gives
\[
 \frac{Q_{l+1}(t)}{tQ_l(t)}
   =-e^{\mathsf T}K^{-2}e
     -t^2 e^{\mathsf T}K^{-4}e
     +O(t^4).
\]
The coefficient of $t^2$ is strictly negative, since
$e^{\mathsf T}K^{-4}e=\|K^{-2}e\|^2>0$. Thus
$Q_{l+1}$ cannot be a constant multiple of $tQ_l$
modulo $t^4$. This proves independence in the odd
two-dimensional space, completing all cases of
\eqref{G15:targetspace}.

Consequently the differential image $\mathcal C_m$
surjects onto $S_m/t^4S_{m-4}$. The determinant-pencil
parameter map followed by third-jet projection has a
surjective differential at the explicit point
\eqref{G15:pencil}. In characteristic zero its image
is dense. This proves third-jet dominance for every
$m\geq2$, and hence for all $m\geq1$.

For $m\geq4$ the third-jet target dimension is
\[
 \sum_{j=0}^3\binom{m-j+2}{2}=2m^2+2.
\]
The fourth-jet target adds
$\binom{m-2}{2}>0$, whereas G-06 bounds
$\dim\mathcal D_{m,4}\leq2m^2+2$. The fourth-jet
map, and therefore every higher full jet map, cannot
be dominant. This establishes the maximal-order
assertion.

Finally apply the coefficient-weight argument of
G-14 with the now proved third-jet dominance. A
degree-$d$ coefficient monomial surviving on
$x_0^a q$ has $x_0$ weight
$ad+j_1+\cdots+j_d$. If $\lambda_1\leq ad+3$,
its evaluation depends only on the first four
coefficient forms $q_0,q_1,q_2,q_3$. Their dense
realization by determinants and block padding force
the same isotypic ideal containment as in G-14.
This proves the multiplicity inequality, including
the elementary cases $m<3$ where all available
jets are included. The rectangular inequality is
the equivalent condition
$nd/4>ad+3$, namely $d(4m-3n)>12$.
\end{proof}

\paragraph{What this resolves.}
The argument proves the needed rank for one explicit
pencil in every size and hence proves dominance. It
does not assume or assert a general strong Lefschetz
theorem for arbitrary Artinian determinantal algebras.
The positivity above is the squared norm for an
ordinary finite-dimensional real symmetric Jacobi
matrix; no arithmetic purity or RH input is used.
No literature-priority claim is made for this deduction.

\begin{proposition}[G-16: Explicit nonrectangular labels throughout the dimension-certified range]\label{G-16}
Let $n>m\geq3$, put $a=n-m$, and suppose
\begin{equation}\label{G16:dimensionhypothesis}
 \dim\mathcal P_{m,4}+3>2n^2+2.
\end{equation}
Define
\[
 b=\binom{n+3}{3}\binom{a+3}{3},\qquad L=n+3a,
 \qquad \beta=6bL^{30},\qquad u=2n^2-10,
\]
and set
\begin{equation}\label{G16:firstdegree}
 C=((n+a)\beta^2)^u,\qquad
 s_*=(u+1)(C-1)+1,\qquad d_* =\beta s_*.
\end{equation}
There is an integer $1\leq d\leq d_*$ with $4\mid md$ such that,
for the explicit shape
\begin{equation}\label{G16:label}
 \lambda_{m,n,d}=
 \left(ad+\frac{md}{4},\frac{md}{4},\frac{md}{4},\frac{md}{4}\right),
\end{equation}
one has
\[
 m_{\lambda_{m,n,d}}(\ell^a\operatorname{per}_m)
       >m_{\lambda_{m,n,d}}(\det_n).
\]
The bound has $\log d_*=O(n^2\log n)$ uniformly for $3\leq m<n$.
There is no additional condition $n<4m/3$ in this proposition.

An explicit progression of successful degrees is also available. Put
\[
 T=\binom{b+\beta}{b}-1,\quad E=4T(\beta!),\quad
 Q=((n+a)E)^{u-1},\quad K=u(Q-1)+1.
\]
For every integer $k\geq K$, the same inequality holds at
$d=Ek$ and the label \eqref{G16:label}.
Because $a>0$, every displayed label is nonrectangular.
\end{proposition}

\paragraph{Scope and sources.}
This applies to every padded pair certified by R-14, R-16, R-18,
R-19, or R-20, including those at or above $n=4m/3$, where G-12
excludes rectangular gaps. The construction uses a Segre--Veronese
coordinate ring and the same Kempf--Ness and reductive-quotient
facts used in G-11. For the quantitative bounds, we use Derksen
(2001), Theorem~1.1 and Proposition~1.2, printed p.~956, and the
image-degree/Chardin argument of G-10 and G-13. The new representation
requires a different coefficient-degree bound, which is checked below.

\begin{proof}
Let $V=\mathbb C^4$, $G=\mathrm{SL}(V)$, and put
\[
 Y_D=\mathcal D_{n,4},\qquad
 Y_P=\{L^a q:L\in V,\ q\in\mathcal P_{m,4}\}^{\rm cl}.
\]
These are the irreducible cones used in R-14, with
\[
 \dim Y_D=2n^2+2,\qquad
 \dim Y_P=\dim\mathcal P_{m,4}+3.
\]
For either cone $Y\subseteq\operatorname{Sym}^n V$, define
\begin{equation}\label{G16:enhancedcone}
 X_Y=\{F\otimes v^a:F\in Y,\ v\in V^*\}^{\rm cl}
 \subseteq W:=\operatorname{Sym}^n V\otimes\operatorname{Sym}^a V^*.
\end{equation}

\paragraph{1. The graded ring and the selected representation.}
The projectivization of \eqref{G16:enhancedcone} is the closed
Segre--Veronese image of $\mathbb P(Y)\times\mathbb P(V^*)$.
In particular $X_Y$ is irreducible and
\begin{equation}\label{G16:xdimension}
 \dim X_Y=\dim Y+3.
\end{equation}
Its coordinate ring has the exact grading
\begin{equation}\label{G16:segrering}
 \mathbb C[X_Y]_d=
       \mathbb C[Y]_d\otimes\operatorname{Sym}^{ad}V.
\end{equation}
Indeed each linear coordinate restricts to a product of a degree-one
coefficient function of $F$ and a degree-$a$ polynomial function of $v$.
Their degree-$d$ products span the right side: the first factor is
generated in degree one, and every degree-$ad$ monomial in $v$ can
be split into $d$ degree-$a$ monomials. The only additional relations
are exactly the restrictions defining this image. This proves
\eqref{G16:segrering} directly, without a projective-normality assumption
on $Y$.

To keep the duals explicit, use the actual coordinate-ring decomposition
\[
 \mathbb C[Y]_d=
       \bigoplus_{\lambda\vdash nd}m_\lambda(Y)S_\lambda V^*.
\]
It is the dual of the positive-weight orbit-span convention used
elsewhere in the paper. Schur's lemma for $G$ and
\eqref{G16:segrering} give
\begin{equation}\label{G16:extraction}
 \dim\mathbb C[X_Y]_d^G=
 \begin{cases}
 m_{\lambda_{m,n,d}}(Y),&4\mid md,\\
 0,&4\nmid md.
 \end{cases}
\end{equation}
For full detail,
\[
 (S_\lambda V^*\otimes\operatorname{Sym}^{ad}V)^G
   =\operatorname{Hom}_G(S_\lambda V,\operatorname{Sym}^{ad}V).
\]
This has dimension one precisely when
$\lambda=(ad+r,r,r,r)$ for an integer $r\geq0$, and otherwise is
zero. The fixed size $|\lambda|=nd$ forces $4r=md$, proving
the formula and ruling out any additional labels in this extraction.

\paragraph{2. Stable points with arbitrary padding.}
Let $x_0,x_1,x_2,x_3$ be an orthonormal basis of $V$ and
$e_0^*$ its first dual vector. Set
\[
 \gamma=\sqrt{\binom{n}{m}},\qquad
 F=x_0^a\bigl(x_0^m+\gamma x_1^m+
                         \gamma x_2^m+\gamma x_3^m\bigr),
 \qquad w=F\otimes(e_0^*)^a.
\]
The degree-$m$ factor of $F$ belongs to both
$\mathcal P_{m,4}$ and $\mathcal D_{m,4}$. In fact each of its two
binary summands factors into $m$ linear factors over $\mathbb C$;
put those factors on the diagonal and the cycle of an $m\times m$
pencil. Its permanent is the sum of the two products, and its
determinant is the same sum after absorbing $(-1)^{m-1}$ into one
cycle factor. Block padding then gives $F\in Y_P\cap Y_D$.
Thus the same tensor $w$ lies in both enhanced cones.

Use the Hermitian structures induced from tensor powers. The four
monomial tensors in $w$ are mutually orthogonal and have equal norm:
$\|x_0^n\|=1$, while
$\|x_0^a x_i^m\|^2=a!m!/n!$ for $i>0$, which is canceled by
$\gamma^2$. Their torus weights are exactly
\[
 ne_0-ae_0=me_0,\qquad
 ae_0+me_i-ae_0=me_i\quad(1\leq i\leq3).
\]
Their mean weight is $(m/4)(1,1,1,1)$, so the traceless diagonal
moment map is zero. The off-diagonal terms vanish as well. On the
first tensor factor, the four monomial weights differ by
$m(e_i-e_j)$, never by one root when $m\geq3$. On the second
factor, an off-diagonal derivative is either zero or orthogonal
to $(e_0^*)^a$. Therefore $w$ has zero $G$ moment map.

Its projective stabilizer has zero Lie algebra. To verify this
without a generic-stabilizer assertion, suppose
$A\in\mathfrak{sl}_4$ and $A\cdot F=\kappa F$.
Write $F=x_0^a q$. Dividing by $x_0^{a-1}$ and reducing modulo
$x_0$ shows that $A\cdot x_0$ is a multiple $\eta x_0$, since
$a>0$ and $q|_{x_0=0}\ne0$. It follows that
$A\cdot q=(\kappa-a\eta)q$.
For the Fermat sum $q$, the monomials $x_j^{m-1}x_i$, $i\ne j$,
are all distinct because $m\geq3$. Their coefficients force every
off-diagonal entry of $A$ to vanish. The four pure-power coefficients
then force every diagonal entry to be the same, and trace zero
forces $A=0$. If $A$ stabilizes the projective tensor line of $w$,
the elementary tangent property of a nonzero decomposable tensor
forces it to stabilize the line of $F$; hence again $A=0$.
Thus the projective stabilizer is finite.

Kempf--Ness now makes $w$ a stable point of each enhanced cone.
Their stable loci are nonempty open sets. Using
\eqref{G16:xdimension}, the two affine invariant quotients therefore
have dimensions
\begin{equation}\label{G16:quotdims}
 \begin{split}
 \dim A_D&=u=2n^2-10,\\
 \dim A_P&=\dim\mathcal P_{m,4}-9>u,
 \end{split}
 \qquad
 A_D=\mathbb C[X_{Y_D}]^G,\quad A_P=\mathbb C[X_{Y_P}]^G.
\end{equation}
The subtraction is $15=\dim G$. The padding weights were canceled
in the tensor construction, so this calculation needs no inequality
between $n$ and $4m/3$.

\paragraph{3. A quantitative invariant bound for this representation.}
The dimension of $W$ is $b$. In the matrix coordinates on
$\mathrm{SL}_4$, the action on $V^*$ is given by the transpose
adjugate, whose entries have degree three. The action on $W$
therefore has matrix coefficients of degree at most
$L=n+3a$. Its kernel is finite: its differential is a nonzero
representation of the simple Lie algebra $\mathfrak{sl}_4$, hence
injective, and a zero-dimensional algebraic kernel is finite.
For example the vector $x_0^n\otimes(e_1^*)^a$ has a nonzero
diagonal Lie-algebra weight, verifying nontriviality.

Derksen's Proposition~1.2 applies to the embedding
$\mathrm{SL}_4\subseteq\mathbb A^{16}$ defined by the degree-four
equation $\det-1=0$. Thus its parameters are $t=16$, group dimension
$15$, equation-degree bound $4$, and action-degree bound $L$, giving
$\sigma(W)\leq4L^{15}$.
Theorem~1.1, with invariant-quotient dimension at most $b$, then gives
generators of $\mathbb C[W]^G$ of degrees at most
\[
 \max\{2,(3/8)b(4L^{15})^2\}\leq6bL^{30}=\beta.
\]
Exactness of invariants for this reductive group makes $A_D$ and
$A_P$ graded quotients of that invariant ring. Choose common
homogeneous generators $g_i$ of degrees $1\leq e_i\leq\beta$.

The determinant enhanced cone has a dense parametrization
\[
 (Z,v)\longmapsto\det(Z_1x_0+Z_2x_1+Z_3x_2+Z_4x_3)\otimes v^a.
\]
Its coordinates are homogeneous of degree $n+a$ in the
$4n^2+4$ parameters. This parameter degree is distinct from the
group-action coefficient degree $n+3a$ above.
The generator coordinates $g_i$ composed with this map have degree
$(n+a)e_i\leq(n+a)\beta$.
Homogenizing them with one additional source coordinate gives a
dominant rational map to the ordinary projective closure of
$\operatorname{Spec}A_D$ in its generator embedding. The image-degree
argument of G-10, including its treatment of a base locus, yields
\[
 \deg\overline{\operatorname{Spec}A_D}
      \leq((n+a)\beta)^u.
\]

Let $H_D^{\leq}(d)=\sum_{j=0}^d\dim(A_D)_j$, and define
$H_P^{\leq}$ similarly. The graded pieces through degree $d$
are contained in the ordinary degree-at-most-$d$ generator filtration,
because all generator weights are at least one. Chardin's bound on
the ordinary projective closure therefore gives
\begin{equation}\label{G16:cumulativeupper}
 H_D^{\leq}(d)\leq((n+a)\beta)^u\binom{d+u}{u}.
\end{equation}
By \eqref{G16:quotdims}, at least $u+1$ images of the generators in
$A_P$ are algebraically independent. Their monomials of ordinary
total degree at most $s$ have graded degree at most $\beta s$;
hence
\begin{equation}\label{G16:cumulativelower}
 H_P^{\leq}(\beta s)\geq\binom{s+u+1}{u+1}.
\end{equation}
As in G-13, the exact binomial comparison is
\[
 \frac{\binom{s+u+1}{u+1}}{\binom{\beta s+u}{u}}
 =\frac{s+u+1}{u+1}\prod_{i=1}^u\frac{s+i}{\beta s+i}
 \geq\frac{s+u+1}{(u+1)\beta^u}.
\]
At $s=s_*$ this is strictly greater than
$C/\beta^u=((n+a)\beta)^u$. Equations
\eqref{G16:cumulativeupper}--\eqref{G16:cumulativelower} give a
strict cumulative gap through $d_*$. Degree zero contributes one
on both sides, so there is a positive degree $d\leq d_*$ with
$\dim(A_P)_d>\dim(A_D)_d$.
Formula \eqref{G16:extraction} identifies this gap with exactly the
label \eqref{G16:label} and forces $4\mid md$.
The four-variable highest-weight restriction of G-05/R-14 transfers
it to the full padded-permanent and determinant orbit closures.
Since $a<n$, $b=O(n^6)$, $L=O(n)$, and $\beta=O(n^{36})$;
\eqref{G16:firstdegree} consequently has
$\log d_*=O(n^2\log n)$.

\paragraph{4. An explicit progression, if individual degrees are desired.}
There are at most $T=\binom{b+\beta}{b}-1$ generators, since the
ambient polynomial space in positive degrees through $\beta$ has
that dimension. The elementary generator-grouping proof in G-11
therefore makes both $E$th Veronese rings standard graded for
$E=4T(\beta!)$. Their dimensions are still $u$ and at least $u+1$.
The degree-one invariant functions of the first Veronese, composed
with the enhanced determinant parametrization, have degree
$(n+a)E$. The projective image has dimension $u-1$ and degree at
most $Q=((n+a)E)^{u-1}$.
Chardin's theorem and linear Noether normalization give respectively
\[
 \dim(A_D)_{Ek}\leq Q\binom{k+u-1}{u-1},\qquad
 \dim(A_P)_{Ek}\geq\binom{k+u}{u}.
\]
Their binomial ratio is $(k+u)/u$, strictly larger than $Q$ for
$k\geq u(Q-1)+1$. This proves the claimed progression; $4\mid E$
ensures all labels in it are integral.
\end{proof}

\paragraph{Examples where ordinary rectangles are excluded.}
The pairs $(m,n)=(15,20)$ and $(16,22)$ from R-17 satisfy
\eqref{G16:dimensionhypothesis}. G-16 produces labels
$\lambda=(35r,15r,15r,15r)$ at degrees $d=4r$ for the first pair,
and labels $\lambda=(10d,4d,4d,4d)$ for the second, at the
successful degrees specified above. G-12 precludes ordinary
rectangular gaps for both pairs. The same explicit label family
applies to the unbounded nonrectangular strip supplied by R-20.


\begin{theorem}[G-17: Exact criterion for four-row stretched directions]\label{G-17}
Assume $n>m\geq3$, put $a=n-m$, and suppose
\[
 \dim\mathcal P_{m,4}+3>2n^2+2.
\]
Fix $s\geq1$ and a partition $\nu\vdash ns$ of length at most four,
appending zero parts when necessary. There is a positive integer $j$
such that
\begin{equation}\label{G17:gap}
 m_{j\nu}(\ell^a\operatorname{per}_m)>m_{j\nu}(\det_n),
 \qquad\text{in coordinate degree }sj,
\end{equation}
if and only if
\begin{equation}\label{G17:criterion}
 \nu_4>0\qquad\hbox{and}\qquad\nu_1>as.
\end{equation}
If either condition fails, no positive stretch $j$ gives a gap.
If both hold, an explicit progression of stretches gives gaps at
every sufficiently large term, as specified below. This does not
assert a gap at every sufficiently large integer stretch.

Here are effective bounds in the affirmative case. Put
\[
 \eta=(\nu_1-\nu_4,\nu_2-\nu_4,\nu_3-\nu_4,0),\qquad e=|\eta|,
 \qquad r_\eta=\#\{(i,j):1\leq i<j\leq4,\ \eta_i>\eta_j\},
\]
and define
\[
 \begin{aligned}
 b&=\binom{\binom{n+3}{3}+s-1}{s}\dim S_\eta\mathbb C^4,
 &u&=2n^2+r_\eta-13,\\
 L&=ns+3e,&p&=ns+e,\qquad \beta=6bL^{30}.
 \end{aligned}
\]
The Schur dimension is explicitly
$\prod_{i<j}(\eta_i-\eta_j+j-i)/(j-i)$.
One successful stretch satisfies
\begin{equation}\label{G17:firststretch}
 1\leq j\leq
 j_*:=\beta\bigl[(u+1)((p\beta^2)^u-1)+1\bigr].
\end{equation}
For a progression, let
\[
 T=\binom{b+\beta}{b}-1,\quad E=4T(\beta!),\quad
 Q=(pE)^{u-1},\quad K=u(Q-1)+1.
\]
Then \eqref{G17:gap} holds for every stretch $j=Ek$ with $k\geq K$.
\end{theorem}

\paragraph{Meaning of the criterion.}
For any one of the dimension-certified padded pairs, this determines
exactly which rational directions in the four-row partition simplex
contain a stretched multiplicity obstruction. Its two excluded faces
come from low-row determinant survival and the padding weight bound.
In particular this is a classification of all such directions,
not only a further special choice of the one-longer-row labels in G-16.
The coordinate degrees and stretch factors remain large; no efficient
multiplicity-positivity algorithm is asserted.

\begin{proof}
Necessity follows from earlier ideal containments and includes equality.
If $\nu_4=0$, every $j\nu$ has at most three rows, and G-05 gives
the full ambient multiplicity on the determinant side. If
$\nu_1\leq as$, then in coordinate degree $d=sj$ one has
$(j\nu)_1\leq ad$. The coefficient-weight argument of G-14 applies
and gives the opposite weak multiplicity inequality for every $j$.
Thus it remains to prove sufficiency under \eqref{G17:criterion}.

Use $V=\mathbb C^4$ and the two degree-$n$ cones $Y_D,Y_P$ of G-16.
We first give the auxiliary flag construction, including its coordinate
ring, so that no projective-normality or multiplicity-free assertion
is left implicit.

\paragraph{1. The auxiliary cone and its Cartan coordinate ring.}
Let $v_\eta$ be a nonzero extremal vector of weight $-\eta$ in
$S_\eta V^*$. Concretely, put
$c_i=\eta_i-\eta_{i+1}$ for $i=1,2,3$ and take the product of
the vectors
\[
 (e_0^*\wedge\cdots\wedge e_{i-1}^*)^{\otimes c_i}.
\]
It lies in the single Cartan component $S_\eta V^*$ of the indicated
tensor product of exterior powers. Its projective orbit records the
nested subspaces of dimensions $i$ for which $c_i>0$.
This is the corresponding partial flag variety, a closed incidence
subvariety of a product of Grassmannians. Its dimension is
$r_\eta$: its stabilizer is the block parabolic for the equal-part
blocks of $\eta$, and the missing off-diagonal blocks contribute
one parameter for each pair counted in $r_\eta$.
If $\eta=0$, the projective orbit is a point and $r_\eta=0$.

Let $\mathcal C_\eta$ be the affine cone over this projective orbit.
Its degree-$d$ coordinate ring is
\begin{equation}\label{G17:cartanring}
 \mathbb C[\mathcal C_\eta]_d=S_{d\eta}V.
\end{equation}
Here is a direct representation-theoretic proof. The dual of its
degree-$d$ coordinate space is the span of the group translates of
$v_\eta^{\otimes d}$. This is an extremal vector of weight $-d\eta$.
Its weight space at the extreme of the tensor-product weight order
is one dimensional, and it generates exactly the irreducible Cartan
component $S_{d\eta}V^*$. Taking the dual proves
\eqref{G17:cartanring}. Products give its usual Cartan multiplication.
When $\eta=0$, the cone is a line and each graded piece is
one dimensional, so the same formula holds.

For either cone $Y$, define the enhanced Veronese cone
\begin{equation}\label{G17:enhancement}
 X_Y=\{F^{\otimes s}\otimes v:F\in Y,\ v\in\mathcal C_\eta\}^{\rm cl}
 \subseteq W:=\operatorname{Sym}^s(\operatorname{Sym}^n V)
                      \otimes S_\eta V^*.
\end{equation}
Its projectivization is the Segre product of the $s$th Veronese
embedding of $\mathbb P(Y)$ with the partial flag variety. Hence
\begin{equation}\label{G17:dimension}
 \dim X_Y=\dim Y+r_\eta.
\end{equation}
Pullback of degree-$d$ coordinate functions gives exactly
\begin{equation}\label{G17:ring}
 \mathbb C[X_Y]_d=\mathbb C[Y]_{sd}\otimes S_{d\eta}V.
\end{equation}
The first factor is the Veronese subring of the standard graded
coordinate ring of $Y$, and the second is \eqref{G17:cartanring}.
Both are generated by their degree-one pieces in these gradings;
products of paired coordinates span their diagonal tensor product.
The pullback into the product coordinate ring is injective, proving
\eqref{G17:ring} without assuming that $Y$ is normal.

Write the actual degree-$sd$ coordinate representations as
$\bigoplus_{\lambda\vdash nsd}m_\lambda(Y)S_\lambda V^*$.
The invariant multiplicity in \eqref{G17:ring} is
\[
 \dim\operatorname{Hom}_{\mathrm{SL}_4}
                   (S_\lambda V,S_{d\eta}V).
\]
It is one if $\lambda=d\eta+t(1,1,1,1)$, and zero otherwise.
The size $nsd$ forces
$4t=(ns-e)d=4\nu_4d$, so $\lambda=d\nu$.
Consequently
\begin{equation}\label{G17:extraction}
 \dim\mathbb C[X_Y]_d^{\mathrm{SL}_4}=m_{d\nu}(Y).
\end{equation}
This is a single-label extraction, with original coordinate degree
$sd$, rather than a sum over different labels or degrees.

\paragraph{2. Stability for every strict direction in the criterion.}
Define four positive rational numbers
\[
 \theta_0=\frac{\nu_1/s-a}{m},\qquad
 \theta_i=\frac{\nu_{i+1}}{sm}\quad(1\leq i\leq3).
\]
Positivity is exactly \eqref{G17:criterion}, and their sum is one.
Let $u_0=x_0^n$ and
$u_i=\sqrt{\binom nm}\,x_0^a x_i^m$ for $i>0$.
These are orthonormal monomial vectors for the tensor-induced metric.
Put
\[
 F=\sum_{i=0}^3\sqrt{\theta_i}\,u_i,
 \qquad w=F^{\otimes s}\otimes v_\eta,
\]
normalizing $v_\eta$ to have norm one. The degree-$m$ factor of $F$
is a Fermat sum with four nonzero coefficients. The binary factorization
and cyclic-pencil construction in G-16 therefore put $F$ in both
$Y_D$ and $Y_P$, so $w$ belongs to both enhanced cones.

The diagonal moment of $F$, before removing the scalar part, is
\[
 a e_0+m\sum_i\theta_i e_i=\nu/s.
\]
The moment of a unit pure tensor power is $s$ times the moment of
the original vector; this follows by applying the Lie operator to
each factor and taking the tensor inner product. Thus the first
factor of $w$ contributes $\nu$. The auxiliary extremal vector
contributes its weight $-\eta$. Their sum is
$\nu-\eta=\nu_4(1,1,1,1)$, whose traceless part is zero.
Off-diagonal moments vanish on $F$ because its monomial weights
differ by $m$ times a root, with $m\geq3$, and on $v_\eta$ because
it is a weight vector. They therefore vanish on the tensor product.

The projective stabilizer of this $F$ has zero Lie algebra by the
mixed-monomial calculation in G-16; all four coefficients remain
nonzero, as required there. A Lie operator stabilizing the line of
$F^{\otimes s}$ must stabilize the line of $F$: in the symmetric
algebra on $\operatorname{Sym}^n V$ the equation
$sF^{s-1}(A\cdot F)=cF^s$ implies $A\cdot F=(c/s)F$.
A projective stabilizer of the decomposable tensor $w$ consequently
stabilizes that same line. Thus its projective stabilizer is finite.
Kempf--Ness gives a stable point of each enhanced cone.

The invariant quotients have dimensions
\begin{equation}\label{G17:quotdims}
 \begin{split}
 \dim A_D&=2n^2+r_\eta-13=u,\\
 \dim A_P&=\dim\mathcal P_{m,4}+r_\eta-12>u,
 \end{split}
 \qquad A_D=\mathbb C[X_{Y_D}]^{\mathrm{SL}_4},\quad
 A_P=\mathbb C[X_{Y_P}]^{\mathrm{SL}_4}.
\end{equation}
Indeed stable loci are nonempty and open, so one subtracts the
common orbit dimension $15$ from \eqref{G17:dimension}.
The original dimension gap is preserved for every flag type.

\paragraph{3. The two polynomial degrees and the effective bounds.}
The dimension of $W$ is $b$. Its group-action matrix coefficients
have degree at most $L=ns+3e$ in the entries of $g\in\mathrm{SL}_4$:
the first factor has degree $ns$, while $S_\eta(g^{-\mathsf T})$
has degree at most $3e$ because each entry of $g^{-\mathsf T}$ is
a cubic cofactor. The representation has finite kernel, since its
kernel is contained in the finite projective stabilizer of $w$.
Derksen's Theorem~1.1 and Proposition~1.2 therefore give generators
of $\mathbb C[W]^{\mathrm{SL}_4}$ of degrees at most
$\beta=6bL^{30}$, exactly as in G-16. Exactness of invariants makes
$A_D,A_P$ graded quotients with the same bound on their generators.

For the image-degree bound a smaller, different parameter degree
is valid, and the reason is important. Introduce an independent
matrix $B\in\operatorname{End}(V^*)$, without imposing any relation
between $B$ and the determinant-pencil parameters. The map
\[
 B\longmapsto S_\eta(B)v_\eta
\]
is a homogeneous polynomial map of degree $e$. One can see this
directly in the tensor construction of the Schur functor: it is
the restriction of $B^{\otimes e}$ to the relevant Schur component.
The invertible $B$ form a dense open set and give the full
$\mathrm{GL}(V^*)$ orbit, which is dense in $\mathcal C_\eta$
when $e>0$. Its scalar matrices supply every nonzero scalar
multiple, because the base field is $\mathbb C$.
Thus
\begin{equation}\label{G17:parameterization}
 (Z,B)\longmapsto
  \det(Z_0x_0+\cdots+Z_3x_3)^{\otimes s}
       \otimes S_\eta(B)v_\eta
\end{equation}
has dense image in $X_{Y_D}$ and is homogeneous of degree
$p=ns+e$ in its $4n^2+16$ independent parameters.
No inverse-matrix constraint or adjugate is needed for this
parametrization. If $e=0$, the second factor is constant and can
be omitted; all scalar multiples are already supplied by scaling
the first cone. Formula \eqref{G17:parameterization} then remains
valid with degree $p=ns$.

Apply the cumulative generator-filtration comparison from G-16,
now with quotient dimension $u$, generator bound $\beta$, and
parameter degree $p$. For clarity, its two inequalities become
\[
 \sum_{d=0}^{D}\dim(A_D)_d
       \leq(p\beta)^u\binom{D+u}{u},\qquad
 \sum_{d=0}^{\beta h}\dim(A_P)_d
       \geq\binom{h+u+1}{u+1}.
\]
The first uses the ordinary projective closure in the generator
embedding, of degree at most $(p\beta)^u$, and Chardin's theorem.
The second uses $u+1$ algebraically independent generator images.
The same exact binomial ratio is at least
$(h+u+1)/((u+1)\beta^u)$. Taking
$h=(u+1)((p\beta^2)^u-1)+1$ yields a strict cumulative gap through
$j_*=\beta h$. Since degree zero agrees, some positive stretch
$j\leq j_*$ has a strict invariant-dimension gap. Equation
\eqref{G17:extraction} proves \eqref{G17:gap} at that stretch.

For the individual progression, there are at most
$T=\binom{b+\beta}{b}-1$ homogeneous invariant generators.
The explicit grouping argument of G-11 makes the $E$th Veronese
rings standard graded for $E=4T(\beta!)$.
Their dimensions are $u$ and at least $u+1$, and composing their
degree-one invariant coordinates with \eqref{G17:parameterization}
gives parameter degree $pE$. Consequently their degree-$k$ Hilbert
functions satisfy
\[
 \dim(A_D)_{Ek}\leq Q\binom{k+u-1}{u-1},\qquad
 \dim(A_P)_{Ek}\geq\binom{k+u}{u},\qquad Q=(pE)^{u-1}.
\]
The ratio is strictly greater than one once
$k\geq u(Q-1)+1$. Applying \eqref{G17:extraction} gives the
specified successful multiples $j=Ek$.
On this progression the determinant-side multiplicity is positive as
well. The positive-dimensional graded domain $A_D$ has a nonzero
positive-degree generator $g$ of some degree $e_0\leq\beta$.
Since $e_0\mid E$, the elements $g^{Ek/e_0}$ are nonzero for every
$k\geq1$. Equation \eqref{G17:extraction} therefore gives
$m_{Ek\nu}(\det_n)\geq1$. Thus the displayed strict gaps compare
two positive multiplicities and are not occurrence gaps.
Finally the four-variable highest-weight restriction transfers
all the multiplicities back to the full orbit closures.
\end{proof}

\paragraph{Consistency with the earlier boundaries.}
If $\nu$ is rectangular, the strict inequality in
\eqref{G17:criterion} is $n/4>a$, equivalently $n<4m/3$,
recovering G-11/G-12 whenever the dimension hypothesis holds.
For the one-longer-row direction in G-16, it holds for all padding
because the extra first-row component compensates for $a$.
At equality $\nu_1=as$, every stretch remains excluded, rather
than merely failing this stability construction.
G-15 further excludes a finite initial segment when
$j(\nu_1-as)\leq3$; this is consistent with the eventual
progression in the strict interior.

